\subsection{PRIMITIVES OF DIFFERENTIAL FORMS WITH VALUES IN A LIE ALGEBRA}

Lemma 5. Let X be a manifold of class \(\mathbf{C}^r\) , F and \(\mathbf{F}'\) vector bundles of class \(\mathbf{C}^r\) with base space X and \(\phi\) a morphism of F into \(\mathbf{F}'\) . For all \(x \in \mathbf{X}\) , let \(\mathbf{S}_x\) be the set of

\[
(a, \phi (a)) \in \mathrm{F} _ {x} \oplus \mathrm{F} _ {x} ^ {\prime}
\]

for \(a \in \mathbf{F}_x\) . Then the union \(\mathbf{S}\) of the \(\mathbf{S}_x\) is a vector subbundle of \(\mathbf{F} \oplus \mathbf{F}'\) .

Let \(\theta\) and \(\theta'\) be the mappings of \(F \oplus F'\) into itself defined as follows: if \((u, v) \in F_x \oplus F_x'\) , then

\[
\theta (u, v) = (u, v + \phi (u)), \quad \theta^ {\prime} (u, v) = (u, v - \phi (u)).
\]

By Differentiable and Analytic Manifolds, R, 7.7.1, \(\theta\) and \(\theta'\) are morphisms of \(F \oplus F'\) into itself. Clearly \(\theta \circ \theta' = \theta' \circ \theta = \mathrm{Id}_{F\oplus F'}\) . Hence \(\theta\) and \(\theta'\) are automorphisms of \(F \oplus F'\) . Therefore, \(S = \theta(F \oplus \{0\})\) is a vector subbundle of \(F \oplus F'\) .

Lemma 6. Let G be a Lie group germ, \(\omega\) the canonical left differential form of G (§ 3, no. 18.9), M a manifold of class \(\mathbf{C}^r\) ( \(r \geqslant 2\) ) and \(\alpha\) a differential form of class \(\mathbf{C}^{r-1}\) and degree 1 on M with values in L(G).

\begin{enumerate}
\def\labelenumi{(\roman{enumi})}
\tightlist
\item
  The elements of \(\mathbf{T}(\mathbf{M}\times \mathbf{G})\) at which the differential form
\end{enumerate}

\[
\theta = \mathrm{pr} _ {1} ^ {*} \alpha - \mathrm{pr} _ {2} ^ {*} \omega
\]

is zero constitute a vector subbundle S of T(M × G) of class \(\mathbf{C}^{r-1}\) .

\begin{enumerate}
\def\labelenumi{(\roman{enumi})}
\setcounter{enumi}{1}
\item
  For all \((x,g)\in \mathbf{M}\times \mathbf{G},\mathrm{T}(\mathrm{pr}_1)|\mathrm{S}_{(x,g)}\) is an isomorphism of \(\mathrm{S}_{(x,g)}\) onto \(\mathrm{T}_x(\mathbf{M})\) .
\item
  If \(d\alpha + [\alpha]^2 = 0\) (cf.~§3, no. 14) the vector subbundle S is integrable.
\end{enumerate}

If \((x,g)\in \mathbf{M}\times \mathbf{G}\) and \((u,v)\in \mathrm{T}_x(\mathrm{M})\times \mathrm{T}_g(\mathrm{G})\) , then

\[
\theta_ {(x, g)} (u, v) = \alpha (u) - g ^ {- 1} v.
\]

Hence the kernel of \(\theta_{(x,g)}\) is the set \(\mathrm{S}_{(x,g)}\) of \((u,g\alpha (u))\) for \(u\in \mathrm{T}_x(\mathbf{M})\) , whence (ii). We consider \(\mathrm{T}(\mathrm{M}\times \mathrm{G})\) as the direct sum of two vector bundles F and \(\mathbf{F}'\) with \(\mathrm{F}_{(x,g)} = \mathrm{T}_x(\mathrm{M})\times \{0\}\) and \(\mathrm{F}_{(x,g)}^{\prime} = \{0\} \times \mathrm{T}_g(\mathrm{G})\) for all

\[
(x, g) \in \mathbf {M} \times \mathbf {G}.
\]

For \(u \in \mathrm{T}_{x}(\mathbf{M}) \times \{0\}\) , we write \(\phi(u) = (0, g\alpha(u))\) . Using the left trivialization of \(\mathrm{T}(\mathrm{G})\) , it is seen that \(\phi\) is a morphism of \(\mathbf{F}\) into \(\mathbf{F}'\) , whence (i) (Lemma 5). Finally, if \(d\alpha + [\alpha]^2 = 0\) , then

\[
\begin{array}{r l} {d \theta = \mathrm{pr} _ {1} ^ {*} (d \alpha) - \mathrm{pr} _ {2} ^ {*} (d \omega)} \\ & {= - \frac {1}{2} (\mathrm{pr} _ {1} ^ {*} \alpha \wedge \mathrm{pr} _ {1} ^ {*} \alpha - \mathrm{pr} _ {2} ^ {*} \omega \wedge \mathrm{pr} _ {2} ^ {*} \omega)} \\ & {= - \frac {1}{2} (\mathrm{pr} _ {1} ^ {*} \alpha - \mathrm{pr} _ {2} ^ {*} \omega) \wedge (\mathrm{pr} _ {1} ^ {*} \alpha + \mathrm{pr} _ {2} ^ {*} \omega)} \\ & {= - \frac {1}{2} \theta \wedge (\mathrm{pr} _ {1} ^ {*} \alpha + \mathrm{pr} _ {2} ^ {*} \omega)} \end{array}
\]

and hence S is integrable (Differentiable and Analytic Manifolds, R, 9.3.6).

THEOREM 5. Let \(G\) be a Lie group germ, \(M\) a manifold of class \(C^r (r\geqslant 2)\) and \(\alpha\) a differential form of class \(\mathbf{C}^{r-1}\) and degree 1 on \(M\) with values in \(L(G)\) , such that \(d\alpha +[\alpha ]^2 = 0\) . For all \(x\in M\) and all \(g\in G\) , there exists a mapping \(f\) , defined and of class \(\mathbf{C}^{r-1}\) on an open neighbourhood of \(x\) , with values in \(G\) , such that \(f(x) = g\) and \(f^{-1}.df = \alpha\) . Two mappings which satisfy these conditions coincide on a neighbourhood of \(x\) .

Let \(x \in M\) and \(g \in G\) . By Lemma 6 (whose notation we adopt) and Differentiable and Analytic Manifolds, R, 9.3.7, there exist an open neighbourhood U of \(x\) in M and a mapping \(m \mapsto \phi(m) = (m, f(m))\) of class \(C^{r-1}\) of U into \(M \times G\) such that \(f(x) = g\) and \(\phi^*(\theta) = 0\) . Then

\[
\begin{array}{l l} f ^ {- 1}. d f = f ^ {*} (\omega) & \text {(§ 3, no. 18.9)} \\ = (\mathrm{pr} _ {2} \circ \phi) ^ {*} (\omega) & \text {(for \(f = \mathrm{pr} _ {2} \circ \phi\))} \\ = \phi^ {*} (\mathrm{pr} _ {1} ^ {*} \alpha - \theta) & \text {(Lemma 6)} \\ = \phi^ {*} (\mathrm{pr} _ {1} ^ {*} \alpha) & \text {(for \(\phi^ {*} (\theta) = 0\))} \\ = \alpha & \text {(for \(\mathrm{pr} _ {1} \circ \phi = \mathrm{Id} _ {\mathbf {U}}\))}. \end{array}
\]

Let \(f'\) be a mapping of class \(\mathbf{C}^{r-1}\) of U into G such that \(f'(x) = g\) and \(f'^{-1}df' = \alpha\) . By § 3, 18.9, \(ff'^{-1}\) is locally constant and hence \(f' = f\) in a neighbourhood of \(x\) .

PROPOSITION 10. Let M be an analytic manifold, g a complete normable Lie algebra and \(\alpha\) an analytic differential form of degree 1 on M, with values in g, with the following properties:

\begin{enumerate}
\def\labelenumi{(\alph{enumi})}
\item
  for all \(m \in \mathbf{M}\) , \(\alpha_{m}\) is an isomorphism of \(\mathrm{T}_{m}(\mathrm{M})\) onto \(\mathfrak{g}\) ;
\item
  \(d\alpha + [\alpha]^{2} = 0.\)
\end{enumerate}

Then, for all \(m_0 \in \mathbf{M}\) , there exist an open neighbourhood \(\mathbf{M}'\) of \(m_0\) in \(\mathbf{M}\) and a Lie group germ structure on \(\mathbf{M}'\) , compatible with the manifold structure on \(\mathbf{M}'\) , with identity element \(m_0\) and with the following properties:

\begin{enumerate}
\def\labelenumi{(\roman{enumi})}
\item
  \(\alpha_{m_{0}}\) is an isomorphism of L(M') onto g;
\item
  the differential form \(m \mapsto \alpha_{m_0}^{-1} \circ \alpha_m\) is the left canonical differential form of \(M'\) . If \(M_1'\) and \(M_2'\) are two such group germs, \(M_1'\) and \(M_2'\) have a common open subgroup germ.
\end{enumerate}

There exists a Lie group germ G such that \(L(G) = g\) . Let \(m_{0} \in M\) . By Theorem 5, there exist an open neighbourhood \(M'\) of \(m_{0}\) in M and an analytic mapping f of \(M'\) into G such that \(f(m_{0}) = e\) and \(f^{-1}.df = \alpha\) . Then \(T_{m_{0}}(f) = \alpha_{m_{0}}\) is an isomorphism of \(T_{m_{0}}(M)\) onto g; hence, shrinking \(M'\) and G, it can be assumed that f is an isomorphism of the manifold \(M'\) onto the manifold G. We transport to \(M'\) the Lie group germ structure on G by means of \(f^{-1}\) . Then \(T_{m_{0}}(f)\) becomes an isomorphism of \(L(M')\) onto \(L(G) = g\) , whence (i). On the other hand, if \(\omega\) denotes the left canonical differential form of G, then

\[
\begin{array}{r} \alpha_ {m _ {0}} ^ {- 1} \circ \alpha_ {m} = (\mathrm{T} _ {m _ {0}} f) ^ {- 1} \circ (f ^ {- 1}. d f) (m) \\ = (\mathrm{T} _ {m} f) ^ {- 1} \circ \omega (f (m)) \circ \mathrm{T} _ {m} f \end{array}
\]

and hence \(m \mapsto \alpha_{m_0}^{-1} \circ \alpha_m\) is the left canonical differential form of \(\mathbf{M}'\) .

Let \(\mathbf{M}^{\prime}\) be an open neighbourhood of \(m_0\) , with a Lie group germ structure, with identity element \(m_0\) and with the analogous properties to properties (i) and (ii). Then \(\alpha_{m_0}\) is an isomorphism of \(\mathrm{L}(\mathbf{M}')\) onto \(\mathfrak{g}\) and also of \(\mathrm{L}(\mathbf{M}'')\) onto \(\mathfrak{g}\) and hence \(\mathrm{L}(\mathbf{M}') = \mathrm{L}(\mathbf{M}'')\) . Therefore, shrinking \(\mathbf{M}'\) and \(\mathbf{M}''\) , it can be assumed that there exists an isomorphism \(\phi\) of the group germ \(\mathbf{M}'\) onto the group germ \(\mathbf{M}''\) (no. 1, Corollary 1 to Theorem 1). Then \(\phi^{-1}.d\phi\) is the canonical left differential of \(\mathbf{M}'\) . On the other hand, let \(\psi\) be the canonical injection of the manifold \(\mathbf{M}' \cap \mathbf{M}''\) into the Lie group germ \(\mathbf{M}''\) ; clearly \(\psi^{-1}.d\psi\) is a restriction of the canonical left differential of \(\mathbf{M}''\) . Hence \((\psi^{-1}.d\psi)(m) = \alpha_{m_0}^{-1} \circ \alpha_m = (\phi^{-1}.d\phi)(m)\) for all \(m \in \mathbf{M}' \cap \mathbf{M}''\) . Therefore \(\phi\) and \(\psi\) coincide on a neighbourhood of \(m_0\) (§ 3, 18.9). This proves the last assertion of the proposition.

COROLLARY. Let M be an analytic manifold of finite dimension n.~Let \(\omega_{1},\ldots,\omega_{n}\) be analytic differential forms of degree 1 on M, with scalar values, which are linearly independent at each point of M and such that, for all \(k=1,\ldots,n\) , \(\omega_{k}\) is a linear combination with constant coefficients of the \(\omega_{i}\wedge\omega_{j}\) . Then, for all \(m_{0}\in M\) , there exists an open neighbourhood \(M'\) of \(m_{0}\) in M and a Lie group germ structure on \(M'\) compatible with the manifold structure on \(M'\) , with identity element \(m_{0}\) and such that \(\omega_{1}|_{M'},\ldots,\omega_{n}|_{M'}\) form a basis of the space of left invariant differential forms on \(M'\) of degree 1 with scalar values.

If \(\mathbf{M}_1'\) and \(\mathbf{M}_2'\) are two such group germs, \(\mathbf{M}_1'\) and \(\mathbf{M}_2'\) have a common open subgroup germ.

Let \(\mathbf{X}_1, \ldots, \mathbf{X}_n\) be the vector fields on M such that, at each point \(m\) of M, the \((\mathbf{X}_i)_m\) constitute the basis of \(\mathrm{T}_m(\mathrm{M})\) dual to \(((\omega_1)_m), \ldots, (\omega_n)_m)\) . These fields are analytic. By hypothesis, there exist \(c_{ijk} \in \mathbf{K}\) ( \(1 \leqslant i,j, k \leqslant n\) ) such that \(c_{ijk} = -c_{jik}\) and \(d\omega_k = \sum_{i < j} c_{ijk}\omega_i \wedge \omega_j\) . By Differentiable and Analytic Manifolds, R, 8.5.7, formula (11),

\[
\langle [ \mathrm{X} _ {i}, \mathrm{X} _ {j} ], \omega_ {k} \rangle = - (d \omega_ {k}) (\mathrm{X} _ {i}, \mathrm{X} _ {j}) = - \left(\sum_ {r <   s} c _ {r s k} \omega_ {r} \wedge \omega_ {s}\right) (\mathrm{X} _ {i}, \mathrm{X} _ {j}) = - c _ {i j k}
\]

and hence \([\mathrm{X}_i,\mathrm{X}_j] = -\sum_{k}c_{ijk}\mathrm{X}_k\) . It then follows that the \(-c_{ijk}\) are the constants of structure of a Lie algebra \(\mathfrak{g}\) relative to a basis \((e_1,\dots ,e_n)\) . For all \(m\in \mathbf{M}\) , let \(\alpha_{m}\) be the linear mapping of \(\mathrm{T}_{m}(\mathrm{M})\) into \(\mathfrak{g}\) which maps \((\mathrm{X}_1)_m\) to \(e_1,\ldots ,(\mathrm{X}_n)_m\) to \(e_n\) . Then \(\alpha\) is an analytic differential form on \(\mathbf{M}\) of degree

1 with values in \(\mathfrak{g}\) and \(\alpha_{m}\) is an isomorphism of \(\mathrm{T}_m(\mathbf{M})\) onto \(\mathfrak{g}\) . On the other hand, \(\alpha = \sum_{k=1}^{n} \omega_k e_k\) and hence

\[
d \alpha = \sum_ {k = 1} ^ {n} (d \omega_ {k}) e _ {k} = \sum_ {k = 1} ^ {n} \left(\sum_ {i <   j} c _ {i j k} \omega_ {i} \wedge \omega_ {j}\right) e _ {k}
\]

and

\[
\begin{array}{r l} {[ \alpha ] ^ {2} = \sum_ {k = 1} ^ {n} [ \omega_ {k} e _ {k} ] ^ {2} + \sum_ {i <   j} (\omega_ {i} e _ {i}) \wedge (\omega_ {j} e _ {j})} & {\text {(§ 3, formula (30))}} \\ {= \sum_ {i <   j} (\omega_ {i} \wedge \omega _ {j}) [ e _ {i}, e _ {j} ]} \\ {= - \sum_ {k = 1} ^ {n} \sum_ {i <   j} (c _ {i j k} \omega_ {i} \wedge \omega_ {j}) e _ {k}} \\ {= - d \alpha.} \end{array}\tag{5}
\]

It then suffices to apply Proposition 10.

\subsection{PASSAGE FROM LAWS OF INFINITESIMAL OPERATION TO LAWS OF OPERATION}

PROPOSITION 11. Let \(\mathrm{G}_1\) and \(\mathrm{G}_2\) be Lie group germs and \(\mathbf{X}_1\) and \(\mathbf{X}_2\) manifolds of class \(\mathrm{C}^r\) ( \(r \geqslant 2\) ). For \(i = 1, 2\) , let \(\psi_i\) be a law chunk of left operation of class \(\mathrm{C}^r\) of \(\mathrm{G}_i\) on \(\mathbf{X}_i\) and \(\mathrm{D}_i\) the associated law of infinitesimal operation. Let \(\mu: \mathrm{G}_1 \to \mathrm{G}_2\) be a morphism and \(\phi: \mathrm{X}_1 \to \mathrm{X}_2\) a mapping of class \(\mathrm{C}^r\) . Suppose that, for all \(a \in \mathrm{L}(\mathrm{G})\) , the vector fields \((\mathrm{D}_{1})_a\) and \((\mathrm{D}_{2})_{\mathrm{L}(\mu)a}\) are \(\phi\) -related (Differentiable and Analytic Manifolds, R, 8.2.6). Then there exists a neighbourhood \(\Omega\) of \(\{e\} \times \mathbf{X}_1\) in \(\mathrm{G}_1 \times \mathbf{X}_1\) such that \(\phi(\psi_1(g,x)) = \psi_2(\mu(g),\phi(x))\) for all \((g,x) \in \Omega\) .

Let \(p_1: \mathrm{G}_1 \times \mathrm{X}_2 \to \mathrm{G}_1\) , \(p_2: \mathrm{G}_1 \times \mathrm{X}_2 \to \mathrm{X}_2\) be the canonical projections. For all \((g_1, x_2) \in \mathrm{G}_1 \times \mathrm{X}_2\) , let \(f_{g_1, x_2}\) be the mapping \(g_1 a \mapsto (\mathrm{D}_2)_{\mathrm{L}(\mu)a}(x_2)\) of \(\mathrm{T}_{g_1}(\mathrm{G}_1) = g_1\mathrm{L}(\mathrm{G}_1)\) into \(\mathrm{T}_{x_2}(\mathrm{X}_2)\) . The \(f_{g_1, x_2}\) define a morphism of \(p_1^* \mathrm{T}(\mathrm{G}_1)\) into \(p_2^* \mathrm{T}(\mathbf{X}_2)\) .

Let \(x_0 \in X_1\) . There exist an open neighbourhood G of e in \(G_1\) and an open neighbourhood X of \(x_0\) in \(X_1\) such that \(\psi_1(g, x)\) and \(\psi_2(\mu(g), \phi(x))\) are defined for \((g, x) \in G \times X\) . We write, for \((g, x) \in G \times X\) ,

\[
\alpha (g, x) = \phi (\psi_ {1} (g, x)) \in \mathrm{X} _ {2}, \quad \beta (g, x) = \psi_ {2} (\mu (g), \phi (x)) \in \mathrm{X} _ {2}.
\]

If \(\mathbf{G}\) and \(\mathbf{X}\) are sufficiently small, then, for all \((a,g,x)\in \mathrm{L}(\mathrm{G}_1)\times \mathrm{G}\times \mathrm{X},\)

\[
\begin{array}{r l} (\mathrm{T} \alpha) (a g, 0 _ {x}) & = (\mathrm{T} \phi) ((\mathrm{D} _ {1}) _ {a} (\psi_ {1} (g, x))) \\ & = (\mathrm{D} _ {2}) _ {L (\mu) a} (\phi (\psi_ {1} (g, x))) \\ & = (\mathrm{D} _ {2}) _ {L (\mu) a} \alpha (g, x), \end{array}
\]

\[
\begin{array}{r l} (\mathrm{T} \beta) (a g, 0 _ {x}) = & (\mathrm{T} \psi_ {2}) (\mathrm{L} (\mu) a. \mu (g), \phi (x)) \\ = & (\mathrm{D} _ {2}) _ {\mathrm{L} (\mu) a} (\psi_ {2} (\mu (g), \phi (x))) \\ = & (\mathrm{D} _ {2}) _ {\mathrm{L} (\mu) a} \beta (g, x). \end{array}
\]

Hence for \(x \in \mathbf{X}\) , the morphisms \(g \mapsto \alpha(g, x)\) and \(g \mapsto \beta(g, x)\) are integrals of \(f\) ; as

\[
\beta (e, x) = \phi (x) = \alpha (e, x)
\]

for all \(x \in \mathbf{X}\) , it follows from Differentiable and Analytic Manifolds, R, 9.3.7, that \(\alpha\) and \(\beta\) coincide on a neighbourhood of \((e, x_0)\) . Hence the proposition.

COROLLARY. Let G be a Lie group germ and X a manifold of class \(\mathbf{C}^{r}\) . Consider two law chunks of left operation of class \(\mathbf{C}^{r}\) of G on X. Suppose that, for all \(a \in L(G)\) , the corresponding vector field \(D_{a}\) on X is the same for both law chunks. Then these two law chunks coincide on a neighbourhood of \(\{e\} \times X\) .

THEOREM 6. Let G be a Lie group germ, X a manifold of class \(\mathbf{C}^r\) ( \(r \geqslant 2\) ) and \(x_0\) a point of X. Let \(a \mapsto \mathrm{D}_a\) be a law of left infinitesimal operation of class \(\mathbf{C}^{r-1}\) of \(\mathrm{L(G)}\) on X.

\begin{enumerate}
\def\labelenumi{(\roman{enumi})}
\item
  There exists an open neighbourhood \(\mathbf{X}'\) of \(x_0\) in \(\mathbf{X}\) and a law chunk of left operation of class \(\mathbf{C}^{r-1}\) of \(\mathbf{G}\) on \(\mathbf{X}'\) such that the associated law of infinitesimal operation is \(a \mapsto \mathbf{D}_a|\mathbf{X}'\) .
\item
  Let there be two law chunks of left operation of class \(\mathbf{C}^{r-1}\) of \(\mathbf{G}\) on an open neighbourhood \(\mathbf{X}''\) of \(x_0\) ; if they admit \(a \mapsto \mathbf{D}_a|\mathbf{X}''\) as associated law of infinitesimal operation, they coincide on a neighbourhood of \((e, x_0)\) .
\end{enumerate}

Assertion (ii) follows from the Corollary to Proposition 11. We prove (i). For all \((g, x) \in \mathbf{G} \times \mathbf{X}\) and all \(a \in \mathrm{L}(\mathbf{G})\) , we write

\[
\mathrm{Q} _ {a} (g, x) = (a g, \mathrm{D} _ {a} (x)) \in \mathrm{T} _ {g} (\mathrm{G}) \times \mathrm{T} _ {x} (\mathrm{X}).
\]

Let \(S_{(g,x)}\) be the set of \(Q_{a}(g,x)\) for \(a\in L(G)\) . By Lemma 5 of no. 6, the \(S_{(g,x)}\) are the fibres of a vector subbundle S of \(T(G)\times T(X)\) . Let a, b be in \(L(G)\) ; then

\[
\begin{array}{r l} {[ \mathrm{Q} _ {a}, \mathrm{Q} _ {b} ] (g, x) = ([ \mathrm{R} _ {a}, \mathrm{R} _ {b} ] (g), [ \mathrm{D} _ {a}, \mathrm{D} _ {b} ] (x))} \\ & = (- \mathrm{R} _ {[ a, b ]} (g), - \mathrm{D} _ {[ a, b ]} (x)) \\ & = \mathrm{Q} _ {- [ a, b ]} (g, x) \end{array}\tag{§ 3, 18.6}
\]

and hence S is integrable (Differentiable and Analytic Manifolds, R, 9.3.3, (iv)). By Differentiable and Analytic Manifolds, R, 9.3.7, there exist an open neighbourhood \(\mathrm{G}_1\) of \(e\) in G, an open neighbourhood \(\mathbf{X}_1\) of \(x_0\) in X and a mapping \((g,x)\mapsto gx\) of class \(\mathrm{C}^{r - 1}\) of \(\mathrm{G}_1\times \mathrm{X}_1\) into X such that \(ex = x\) for all \(x\in \mathbf{X}_1\) and

\[
(a g) x = \mathrm{D} _ {a} (g x) \quad \text { for } a \in \mathrm{L} (\mathrm{G}), g \in \mathrm{G} _ {1}, x \in \mathrm{X} _ {1}.\tag{6}
\]

In particular

\[
a x = \mathrm{D} _ {a} (x).\tag{7}
\]

Let \(G_{2}\) be an open neighbourhood of e in \(G_{1}\) and \(X_{2}\) an open neighbourhood of \(x_{0}\) in \(X_{1}\) such that \(gg'\) is defined and belongs to \(G_{1}\) for g, \(g'\) in \(G_{2}\) and gx is defined and belongs to \(\mathbf{X}_1\) for \((g,x)\in \mathbf{G}_2\times \mathbf{X}_2\) . Consider the mappings \(\alpha_{1},\alpha_{2}\) of \(\mathrm{G}_2\times (\mathrm{G}_2\times \mathrm{X}_2)\) into X defined by

\[
\alpha_ {1} (g, (h, x)) = g (h x), \quad \alpha_ {2} (g, (h, x)) = (g h) x.
\]

They are of class \(\mathbf{C}^{r - 1}\) . Then

\[
\alpha_ {1} (e, (h, x)) = h x = \alpha_ {2} (e, (h, x)).
\]

On the other hand

\[
\begin{array}{r l} \mathrm{T} (\alpha_ {1}) (a g, 0 _ {(h, x)}) & = (a g) (h x) \\ & = \mathrm{D} _ {a} (g (h x)) \qquad \text {by (6)} \\ & = \mathrm{D} _ {a} (\alpha_ {1} (g, (h, x))), \\ \mathrm{T} (\alpha_ {2}) (a g, 0 _ {(h, x)}) & = (a g h) x \\ & = \mathrm{D} _ {a} ((g h) x) \qquad \text {by (6)} \\ & = \mathrm{D} _ {a} (\alpha_ {2} (g, (h, x))). \end{array}
\]

By Differentiable and Analytic Manifolds, R, 9.3.7, \(\alpha_{1}\) and \(\alpha_{2}\) coincide on a neighbourhood of \((e, (e, x_{0}))\) . Then (i) follows from (7) and Proposition 23 of § 1, no. 11.

COROLLARY 1. Let G be a Lie group germ and X a paracompact manifold of class \(\mathbf{C}^r\) ( \(r \geqslant 2\) ). Let \(a \mapsto \mathrm{D}_a\) be a law of left infinitesimal operation of class \(\mathbf{C}^{r-1}\) of L(G) on X.

\begin{enumerate}
\def\labelenumi{(\roman{enumi})}
\item
  There exists a law chunk of left operation of class \(\mathbf{C}^{r-1}\) of \(\mathbf{G}\) on \(\mathbf{X}\) such that the associated law of infinitesimal operation is \(a \mapsto \mathbf{D}_a\) .
\item
  Two laws of left operation of class \(\mathbf{C}^{r-1}\) of G on X which admit \(a \mapsto \mathbf{D}_a\) as associated law of infinitesimal operation coincide on a neighbourhood of \(\{e\} \times \mathbf{X}\) .
\end{enumerate}

Assertion (ii) follows from the Corollary to Proposition 11. By Theorem 6 (i), there exist an open covering \((\mathbf{X}_i)_{i\in I}\) of \(\mathbf{X}\) and, for all \(i\in I\) , a law chunk of left operation \(\psi_i\) of class \(\mathbf{C}^{r-1}\) of \(\mathbf{G}\) on \(\mathbf{X}_i\) such that the associated law of infinitesimal operation is \(a\mapsto \mathrm{D}_a|\mathbf{X}_i\) . As \(\mathbf{X}\) is paracompact, the covering \((\mathbf{X}_i)_{i\in I}\) can be assumed to be locally finite. For all \((i,j)\in I\times I\) and all \(x\in \mathbf{X}_i\cap \mathbf{X}_j\) , \(\psi_i\) and \(\psi_j\) coincide on a neighbourhood of \((e,x)\) (Corollary to Proposition 11). As \(\mathbf{X}\) is normal, we can apply Proposition 24 of § 1, no. 11, which proves (i).

COROLLARY 2. Let X be a paracompact manifold of class \(\mathbf{C}^r\) ( \(r \geqslant 2\) ) and \(\xi\) a vector field of class \(\mathbf{C}^{r-1}\) on X. There exists a law chunk of operation \(\psi\) of class \(\mathbf{C}^{r-1}\) of K on X such that for all \(x \in \mathbf{X}\) , \(\xi(x)\) is the image under \(t \mapsto (t, x)\) of the tangent vector 1 to K at 0. Two law chunks of operation with the above property coincide on a neighbourhood of \(\{0\} \times \mathbf{X}\) .

This is a special case of Corollary 1.

Remark. Laws of left operation can of course be replaced, throughout this no., by laws of right operation.

\section{FORMAL CALCULATIONS IN LIE GROUPS}

Let \(f, g\) be two formal power series with coefficients in K in the same indeterminates, let \(f_i\) (resp. \(g_i\) ) be the homogeneous component of \(f\) (resp. \(g\) ) of degree i. We shall write

\[
f \equiv g \mod \deg p
\]

\[
\text { if } f _ {i} = g _ {i} \text { for } i <   p.
\]

In this §, G denotes a Lie group germ of finite dimension n; the base field K is assumed to be of characteristic zero. We identify once and for all, by means of a chart, an open neighbourhood of e in G with an open neighbourhood U of 0 in \(K^{n}\) , so that e is identified with 0. For x, y in U and \(n \in \mathbf{Z}\) , x.y denotes the product of x and y and \(x^{[m]}\) the m-th power of x in G (when they are defined). The coordinates of \(x \in U\) are denoted by \(x_{1}, x_{2}, \ldots, x_{n}\) .

\subsection{\texorpdfstring{THE COEFFICIENTS \(c_{\alpha\beta\gamma}\)}{THE COEFFICIENTS c\_\{\textbackslash alpha\textbackslash beta\textbackslash gamma\}}}

Let \(\Omega\) be the set of \((x,y)\in\mathbf{U}\times\mathbf{U}\) such that \(x.y\) is defined and belongs to \(\mathbf{U}\) . Then \(\Omega\) is open in \(\mathbf{U}\times\mathbf{U}\) and the mapping \((x,y)\mapsto x.y\) of \(\Omega\) into \(\mathbf{U}\) is analytic. The coordinates \(z_{1},\ldots,z_{n}\) of \(z=x.y\) therefore admit expansions as integral series about the origin in the powers of \(x_{1},\ldots,x_{n},y_{1},\ldots,y_{n}\) . Therefore, there exist well defined constants \(c_{\alpha_{1},\ldots,\alpha_{n},\beta_{1},\ldots,\beta_{n}}\in\mathbf{K}\) , such that

\[
z _ {1} ^ {\gamma_ {1}} \dots z _ {n} ^ {\gamma_ {n}} = \sum_ {\alpha_ {1},\dots,\alpha_ {n},\beta_ {1},\dots,\beta_ {n} \in \mathbf {N}} c _ {\alpha_ {1},\dots,\alpha_ {n},\beta_ {1},\dots,\beta_ {n},\gamma_ {1},\dots,\gamma_ {n}} x _ {1} ^ {\alpha_ {1}} \dots x _ {n} ^ {\alpha_ {n}} y _ {1} ^ {\beta_ {1}} \dots y _ {n} ^ {\beta_ {n}}\tag{1}
\]

for \(\gamma_{1},\ldots,\gamma_{n}\) in \(\mathbf{N}\). Adopting the conventions of Differentiable and Analytic Manifolds, R, we shall write these formulae more briefly:

\[
(x. y) ^ {\gamma} = \sum_ {\alpha , \beta \in \mathbf {N} ^ {n}} c _ {\alpha , \beta , \gamma} x ^ {\alpha} y ^ {\beta} \quad (\gamma \in \mathbf {N} ^ {n}).\tag{2}
\]

Since \(x.0 = 0.x = x\) for \(x\in \mathbf{U}\)

\[
c _ {\alpha , 0, \gamma} = c _ {0, \alpha , \gamma} = \delta_ {\alpha \gamma}\tag{3}
\]

where \(\delta_{\alpha \gamma}\) is the Kronecker index. In particular, writing henceforth \(k\) instead of \(\varepsilon_k\) for \(k = 1, \ldots, n\) ,

\[
(x. y) _ {k} = x _ {k} + y _ {k} + \sum_ {| \alpha | \geqslant 1, | \beta | \geqslant 1} c _ {\alpha , \beta , k} x ^ {\alpha} y ^ {\beta}.\tag{4}
\]

Writing \(c_{\alpha \beta} = (c_{\alpha \beta 1}, c_{\alpha \beta 2}, \ldots, c_{\alpha \beta n}) \in \mathbf{K}^n\) , it then follows that

\[
x. y = x + y + \sum_ {| \alpha | \geqslant 1, | \beta | \geqslant 1} c _ {\alpha \beta} x ^ {\alpha} y ^ {\beta}.\tag{5}
\]

On the right hand side of (5), we consider the homogeneous component of degree 2:

\[
\mathrm{B} (x, y) = \sum_ {i, j = 1} ^ {n} c _ {i j} x _ {i} y _ {j}
\]

so that \((x,y)\mapsto \mathrm{B}(x,y)\) is a bilinear mapping of \(\mathbf{K}^n\times \mathbf{K}^n\) into \(\mathbf{K}^n\) . Then

\[
x. y \equiv x + y + \mathrm{B} (x, y) \mod \deg 3.\tag{6}
\]

Formula (4) implies on the other hand that

\[
c _ {\alpha , \beta , \gamma} = 0 \quad \text { if } | \alpha | + | \beta | <   | \gamma |\tag{7}
\]

and that the terms of total degree \(|\gamma|\) in the expansion of \(z^\gamma\) are also those of \((x_1 + y_1)^{\gamma_1}(x_2 + y_2)^{\gamma_2} \ldots (x_n + y_n)^{\gamma_n} = \sum_{\alpha + \beta = \gamma} ((\alpha, \beta)) x^\alpha y^\beta\) (cf.~Differentiable and Analytic Manifolds, R, Notation and conventions). Hence:

\begin{enumerate}
\def\labelenumi{(\arabic{enumi})}
\setcounter{enumi}{7}
\tightlist
\item
\end{enumerate}

\[
c _ {\alpha , \beta , \gamma} = 0 \quad \text { if } \quad | \alpha | + | \beta | = | \gamma | \quad \text { but } \quad \alpha + \beta \neq \gamma\tag{8}
\]

\[
c _ {\alpha , \beta , \alpha + \beta} = ((\alpha , \beta)).\tag{9}
\]

The associativity of the product implies the relation

\[
\sum_ {\alpha , \xi} c _ {\alpha \xi \eta} x ^ {\alpha} \left(\sum_ {\beta , \gamma} c _ {\beta \gamma \xi} y ^ {\beta} z ^ {\gamma}\right) = \sum_ {\xi , \gamma} c _ {\xi \gamma \eta} \left(\sum_ {\alpha , \beta} c _ {\alpha \beta \xi} x ^ {\alpha} y ^ {\beta}\right) z ^ {\gamma}
\]

for all \(x, y, z\) sufficiently close to 0, whence

\[
\sum_ {\xi} c _ {\alpha \xi \eta} c _ {\beta \gamma \xi} = \sum_ {\xi} c _ {\xi \gamma \eta} c _ {\alpha \beta \xi} (\alpha , \beta , \gamma , \eta \text { in } \mathbf {N} ^ {n}).\tag{10}
\]

The group germ G admits a commutative open subgroup germ if and only if \(c_{\alpha\beta\gamma} = c_{\beta\alpha\gamma}\) for all \(\alpha, \beta, \gamma\) in \(\mathbf{N}^{n}\) .

\subsection{BRACKET IN THE LIE ALGEBRA}

For \(\alpha \in \mathbf{N}^n\) , let \(e_{\alpha}\) be the point distribution \(\frac{1}{\alpha!} \frac{\partial^{\alpha}}{\partial x^{\alpha}}\) at the origin. In particular, \(e_k = e_{\varepsilon_k} = \frac{\partial}{\partial x_k}\) . The \(e_{\alpha}\) form a basis of the vector space \(\mathrm{U}(G)\) . If \(f\) is an analytic function on an open neighbourhood of 0 in G and \(f(x) = \sum_{\alpha} \lambda_{\alpha} x^{\alpha}\) is its expansion as an integral series about the origin, then

In particular,

\[
\langle e _ {\alpha}, f \rangle = \lambda_ {\alpha}.
\]

\[
\langle e _ {\alpha}, x ^ {\gamma} \rangle = \delta_ {\alpha \gamma}.
\]

Hence

\[
\begin{array}{r l} \langle e _ {\alpha} * e _ {\beta}, x ^ {\gamma} \rangle & = \langle e _ {\alpha} \otimes e _ {\beta}, (x. y) ^ {\gamma} \rangle \\ & = \langle e _ {\alpha} \otimes e _ {\beta}, \sum_ {\alpha^ {\prime}, \beta^ {\prime}} c _ {\alpha^ {\prime} \beta^ {\prime} \gamma} x ^ {\alpha^ {\prime}} y ^ {\beta^ {\prime}} \rangle \\ & = \sum_ {\alpha^ {\prime}, \beta^ {\prime}} c _ {\alpha^ {\prime} \beta^ {\prime} \gamma} \langle e _ {\alpha}, x ^ {\alpha^ {\prime}} \rangle \langle e _ {\beta}, y ^ {\beta^ {\prime}} \rangle = c _ {\alpha \beta \gamma} \end{array}
\]

and hence

\[
e _ {\alpha} * e _ {\beta} = \sum_ {\gamma} c _ {\alpha \beta \gamma} e _ {\gamma}.\tag{11}
\]

(Formula (10) then expresses associativity on U(G).)

In particular, since L(G) is stable under the bracket.

\[
[ e _ {i}, e _ {j} ] = \sum_ {k} (c _ {i j k} - c _ {j i k}) e _ {k}.\tag{12}
\]

The constants of structure of \(\mathbf{L}(\mathbf{G})\) relative to the basis \((e_{1},\ldots,e_{n})\) are therefore the \(c_{ijk}-c_{jik}\) . In other words, canonically identifying \(\mathbf{L}(\mathbf{G})\) with \(K^{n}\) , we obtain:

\[
[ x, y ] = \mathrm{B} (x, y) - \mathrm{B} (y, x).\tag{13}
\]

PROPOSITION 1.

\begin{enumerate}
\def\labelenumi{(\roman{enumi})}
\tightlist
\item
\end{enumerate}

\[
x ^ {[ - 1 ]} \equiv - x + \mathrm{B} (x, x) \mod \deg 3\tag{i}
\]

\[
x. y. x ^ {[ - 1 ]} \equiv y + [ x, y ] \quad \text {   mod   } \deg 3\tag{ii}
\]

\[
y ^ {[ - 1 ]}. x. y \equiv x + [ x, y ] \quad \text {   mod   } \deg 3\tag{iii}
\]

\[
x ^ {[ - 1 ]}. y ^ {[ - 1 ]}. x. y \equiv [ x, y ] \quad \text {   mod   } \deg 3\tag{iv}
\]

\[
x. y. x ^ {[ - 1 ]}. y ^ {[ - 1 ]} \equiv [ x, y ] \quad \text {   mod   } \deg 3\tag{v}.
\]

(In (i), \(x^{[-1]}\) of course represents the expansion of the function \(x \mapsto x^{[-1]}\) as an integral series about the origin; the other formulae can be interpreted analogously.)

Let \(g_{1}\) and \(g_{2}\) be the homogeneous components of \(x^{[-1]}\) of degrees 1 and 2. Then

\[
\begin{array}{r l} 0 = x. x ^ {[ - 1 ]} \\ \equiv x + g _ {1} (x) + \mathrm{B} (x, g _ {1} (x)) & \text {   mod   } \deg 2 \quad \text {(by (6))} \\ \equiv x + g _ {1} (x) & \text {   mod   } \deg 2 \end{array}
\]

and hence \(g_{1}(x) = -x\) . Then

\[
\begin{array}{r l} 0 & = x. x ^ {[ - 1 ]} \\ & \equiv x + (- x + g _ {2} (x)) + B (x, - x + g _ {2} (x)) \mod \deg 3 \\ & \equiv g _ {2} (x) - B (x, x) \mod \deg 3 \end{array}
\]

and hence \(g_{2}(x) = \mathrm{B}(x, x)\) . This proves (i). Then, using (i),

\[
\begin{array}{l l} x. y. x ^ {[ - 1 ]} \equiv (x + y + \mathrm{B} (x, y)) \cdot (- x + \mathrm{B} (x, x)) & \text {mod} \deg 3 \\ \equiv x + y + \mathrm{B} (x, y) + (- x + \mathrm{B} (x, x)) + \mathrm{B} (x + y, - x) & \text {mod} \deg 3 \\ \equiv y + \mathrm{B} (x, y) - \mathrm{B} (y, x) & \text {mod} \deg 3 \\ \equiv y + [ x, y ] \qquad \qquad \qquad \qquad \qquad \qquad \qquad \qquad \qquad \qquad \qquad \qquad \qquad \qquad \qquad \qquad \qquad \qquad \qquad \qquad \qquad \qquad \qquad \qquad \qquad \qquad \qquad \qquad \qquad \qquad \qquad \qquad \qquad \qquad \qquad \qquad \qquad \\ & (\text {by (13)}) \end{array}
\]

whence (ii). The proof of (iii) is analogous. Combining (i) and (iii), we obtain

\[
\begin{array}{r l r} x ^ {[ - 1 ]}. y ^ {[ - 1 ]}. x. y & \equiv (- x + \mathrm{B} (x, x)). (x + [ x, y ]) & \bmod \deg 3 \\ & \equiv - x + \mathrm{B} (x, x) + x + [ x, y ] + \mathrm{B} (- x, x) & \bmod \deg 3 \\ & \equiv [ x, y ] & \bmod \deg 3 \end{array}
\]

whence (iv). The proof of (v) is analogous.

\subsection{POWERS}

Consider \(j\) points of \(\mathbf{G}\) :

\[
\begin{array}{l} x (1) = (x (1) _ {1}, x (1) _ {2}, \ldots , x (1) _ {n}) \\ x (2) = (x (2) _ {1}, x (2) _ {2}, \ldots , x (2) _ {n}) \\ \cdot \cdot \cdot \\ x (j) = (x (j) _ {1}, x (j) _ {2}, \ldots , x (j) _ {n}). \end{array}
\]

The mapping \((x(1), x(2), \ldots, x(j)) \mapsto x(1).x(2) \ldots x(j)\) admits an expansion as an integral series about the origin:

\[
x (1). x (2) \dots x (j) = \sum_ {\alpha (1), \alpha (2), \dots , \alpha (j) \in \mathbf {N} ^ {n}} a _ {\alpha (1), \dots , \alpha (j)} x (1) ^ {\alpha (1)} \dots x (j) ^ {\alpha (j)}\tag{14}
\]

where the \(a_{\alpha(1), \ldots, \alpha(j)}\) are elements of \(\mathbf{K}^n\) . We write, for \(j = 0, 1, 2, \ldots\) ,

\[
\psi_ {j} (x) = \sum_ {\alpha (1) \neq 0, \dots , \alpha (j) \neq 0} a _ {\alpha (1), \dots , \alpha (j)} x ^ {\alpha (1) + \dots + \alpha (j)}\tag{15}
\]

where the right hand side is a convergent integral series in the variable \(x \in K^{n}\) . This series is obtained by suppressing in (14) the terms in which one of the variables \(x(1), \ldots, x(j)\) does not occur explicitly and then writing \(x(1) = x(2) = \cdots = x(j) = x\) .

If \(t \in K\) , all the t-th power mappings of G have the same expansion as an integral series about the origin, since any two of them coincide on a neighbourhood of 0. We denote this expansion as an integral series by \(x^{[t]}\) .

PROPOSITION 2. (i) \(\psi_{j} \equiv 0 \bmod \deg j\) .

\begin{enumerate}
\def\labelenumi{(\roman{enumi})}
\setcounter{enumi}{1}
\tightlist
\item
  If \(t\in \mathbf{K}\) , then
\end{enumerate}

\[
x ^ {[ t ]} = \sum_ {i = 1} ^ {\infty} \binom {t} {i} \psi_ {i} (x)\tag{16}
\]

where the formal power series on the right is meaningful because of (i).

(We write

\[
\binom {t} {i} = \frac {t (t - 1) \dots (t - i + 1)}{i !}
\]

for all \(t\in \mathbf{K}.\)

Assertion (i) is obvious from the definition of the \(\psi_{j}\) .

We prove (ii) for t an integer \(\geqslant0\) . By (14),

\[
x ^ {[ t ]} = \sum_ {\alpha (1),\ldots,\alpha (t)\in(\mathbf{N}^{n})^{t}} a _ {\alpha (1), \dots , \alpha (t)} x ^ {\alpha (1) + \dots + \alpha (t)}.\tag{17}
\]

For \(\alpha = (\alpha(1),\ldots,\alpha(t))\in(\mathbf{N}^{n})^{t}\) , let \(\sigma(\alpha)\) denote the set of \(j\in\{1,2,\ldots,t\}\) such that \(\alpha(j)\neq0\) . If, in the sum (17), we group together the terms for which \(\sigma(\alpha)\) is the same, then

\[
x ^ {[ t ]} = \sum_ {\sigma \subset (1, t)} h _ {t, \sigma} (x)\tag{18}
\]

where

\[
h _ {t, \sigma} (x) = \sum_ {\sigma (\alpha) = \sigma} a _ {\alpha (1), \dots , \alpha (t)} x ^ {\alpha (1) + \dots + \alpha (t)}.\tag{19}
\]

Let \(\sigma = \{j_1, j_2, \ldots, j_q\}\) with \(j_1 < j_2 < \ldots < j_q\) . In (14) (where \(j\) is replaced by \(t\) ), we substitute 0 for \(x(k)\) for \(k \notin \sigma\) ; as 0 is the identity element of G, we obtain the expansion of \(x(j_1).x(j_2)..x(j_q)\) as an integral series about the origin:

\[
x (j _ {1}). x (j _ {2}) \dots x (j _ {q}) = \sum_ {\sigma (\alpha) \subset \sigma} a _ {\alpha (1), \dots , \alpha (t)} x (j _ {1}) ^ {\alpha (j _ {1})} x (j _ {2}) ^ {\alpha (j _ {2})} \dots x (j _ {q}) ^ {\alpha (j _ {q})}
\]

and hence, by the definition of \(\psi_q\) :

\[
\psi_ {q} (x) = \sum_ {\sigma (\alpha) = \sigma} a _ {\alpha (1), \dots , \alpha (t)} x ^ {\alpha (j _ {1}) + \dots + \alpha (j _ {q})}.\tag{20}
\]

By (19) and (20), we see that \(h_{t,\sigma}(x) = \psi_{\mathrm{Card}\sigma}(x)\) . Then (18) implies

\[
x ^ {[ t ]} = \sum_ {i = 0} ^ {t} \binom {t} {i} \psi_ {i} (x) = \sum_ {i = 0} ^ {\infty} \binom {t} {i} \psi_ {i} (x).
\]

Then we write \(x^{[t]} = \sum_{i=0}^{\infty} \binom{t}{i}\psi_i(x)\) for all \(t\in\mathbf{K}\). In the integral series \(x^{[t]}\) and \(x^{[t']}\), each coefficient is a polynomial function of \(t\). For this is obvious for \(x^{[t']}\). As far as \(x^{[t]}\) is concerned, it suffices to prove that, for all \(u\in\mathrm{U}(G)\), the image of \(u\) under \(x\mapsto x^{[t]}\) is a polynomial function of \(t\). Now, for \(u\in\mathrm{U}^{m}(G)\), this image is \(t^{m}u\) (§ 4, no. 3, Proposition 7 (iv)).

As \(x^{[t]} = x^{[t']}\) for \(t\) an integer \(\geqslant 0\) , it then follows that \(x^{[t]} = x^{[t']}\) for all \(t \in \mathbf{K}\) .

Remarks. (1) We write condition (ii) of Proposition 2 for \(t\) an integer \(\geqslant 0\) :

\[
\begin{array}{r l} & 0 = \psi_ {0} (x) \\ & x = \psi_ {0} (x) + \psi_ {1} (x) \\ & x ^ {[ 2 ]} = \psi_ {0} (x) + 2 \psi_ {1} (x) + \psi_ {2} (x) \\ & \cdot \cdot \cdot . \end{array}
\]

These formulae suffice to determine the \(\psi_{i}\) .

\begin{enumerate}
\def\labelenumi{(\arabic{enumi})}
\setcounter{enumi}{1}
\tightlist
\item
  We see that \(\psi_0(x) = 0\) , \(\psi_1(x) = x\) , \(\psi_2(x) = x^{[2]} - 2x\) ,
\end{enumerate}

\[
x ^ {[ - 1 ]} = \sum_ {i = 1} ^ {\infty} (- 1) ^ {i} \psi_ {i} (x).
\]

\begin{enumerate}
\def\labelenumi{(\arabic{enumi})}
\setcounter{enumi}{2}
\tightlist
\item
  The above expression for \(\psi_{2}\) and formula (6) prove that
\end{enumerate}

\[
\psi_ {2} (x) \equiv \mathrm{B} (x, x) \mod \deg 3.\tag{21}
\]

Using Proposition 2, (i) and (ii), we see that

\[
x ^ {[ t ]} \equiv t x + \binom {t} {2} \mathrm{B} (x, x) \mod \deg 3.\tag{22}
\]

\begin{enumerate}
\def\labelenumi{(\arabic{enumi})}
\setcounter{enumi}{3}
\tightlist
\item
  Let \(\psi_{p,m}(x)\) and \(h_{t,m}(x)\) denote the homogeneous components of \(\psi_p(x)\) and \(x^{[t]}\) of degree \(m\) . Then \(\psi_{p,m} = 0\) for \(m < p\) . On the other hand, Proposition 2 (ii) gives
\end{enumerate}

\[
h _ {t, m} (x) = \sum_ {p \leqslant m} \frac {t (t - 1) \dots (t - p + 1)}{p !} \psi_ {p, m} (x)\tag{23}
\]

that is

\[
h _ {t, m} (x) = \sum_ {1 \leqslant r \leqslant m} t ^ {r} \phi_ {r, m} (x)\tag{24}
\]

where the \(\phi_{r,m}\) are homogeneous polynomial mappings of degree m of \(K^{n}\) into \(K^{n}\) . In particular, by (23),

\begin{enumerate}
\def\labelenumi{(\arabic{enumi})}
\setcounter{enumi}{24}
\tightlist
\item
\end{enumerate}

\[
\phi_ {1, m} (x) = \sum_ {p \leqslant m} \frac {(- 1) ^ {p - 1}}{p} \psi_ {p, m} (x)\tag{25}
\]

\[
\phi_ {m, m} (x) = \frac {1}{m !} \psi_ {m, m} (x).\tag{26}
\]

\begin{enumerate}
\def\labelenumi{(\arabic{enumi})}
\setcounter{enumi}{4}
\tightlist
\item
  If \(\mathbf{K}\) is of characteristic \(>0\) , the results of nos. 1 and 2 remain true, provided \(e_{\alpha}\) , in no. 2, is defined by \(\langle e_{\alpha}, \sum_{\beta} \lambda_{\beta} x^{\beta} \rangle = \lambda_{\alpha}\) . In no. 3, if the \(\psi_j\) are defined as above, the argument again proves that \(x^{[t]} = \sum_{i=1}^{\infty} \binom{t}{i} \psi_i(x)\) if \(t \in \mathbf{N}\) .
\end{enumerate}

\subsection{EXPONENTIAL}

Let \(\mathrm{E}(x)\) be the expansion of an exponential mapping of G as an integral series about 0. Let \(\mathrm{L}(x)\) be the expansion of the inverse mapping of an injective exponential mapping of G as an integral series about 0. Since the tangent mapping at 0 to any exponential mapping is the identity of \(\mathrm{L}(\mathrm{G})\) , \(\mathrm{E}(x) \equiv x \mod \deg 2\) and \(\mathrm{L}(x) \equiv x \mod \deg 2\) . Since \(\mathrm{E}(\mathrm{L}(x)) = \mathrm{L}(\mathrm{E}(x))\) for \(x\) sufficiently close to 0, the formal power series E and L are such that \(\mathrm{E}(\mathrm{L}(\mathrm{X})) = \mathrm{L}(\mathrm{E}(\mathrm{X})) = \mathrm{X}\) . An analogous argument shows that

\[
\mathrm{E} (t \mathrm{X}) = (\mathrm{E} (\mathrm{X})) ^ {[ t ]}, \quad \mathrm{L} (\mathrm{X} ^ {[ t ]}) = t \mathrm{L} (\mathrm{X})
\]

for \(t \in K\) .

PROPOSITION 3.

\begin{enumerate}
\def\labelenumi{(\arabic{enumi})}
\setcounter{enumi}{26}
\tightlist
\item
\end{enumerate}

\[
\mathrm{L} = \sum_ {p = 1} ^ {\infty} \frac {(- 1) ^ {p - 1}}{p} \psi_ {p}\tag{27}
\]

\[
\mathrm{E} = \sum_ {p = 1} ^ {\infty} \frac {1}{p !} \psi_ {p, p}\tag{28}
\]

(Recall that \(\psi_{p,p}\) is the homogeneous component of \(\psi_{p}\) of degree \(p\) .)

\[
\operatorname{E} (t x) = (\operatorname{E} (x)) ^ {[ t ]}
\]

or, by (24),

\[
\mathrm{E} (t x) = \sum_ {m \geqslant 0} \sum_ {1 \leqslant r \leqslant m} t ^ {r} \phi_ {r, m} (\mathrm{E} (x)).
\]

The two sides are formal power series in \(t\) and \(x\) . Equating the terms of first degree in \(t\) , we obtain

\[
x = \sum_ {m \geqslant 0} \phi_ {1, m} (\mathrm{E} (x)).\tag{29}
\]

Now the inversion of a system of formal power series, when it is possible, is possible in only one way (Algebra, Chapter IV, § 6, Corollary to Proposition 8). Then

\[
\begin{array}{r l} \mathrm{L} (x) & = \sum_ {m \geq 0} \phi_ {1, m} (x) \quad \text { by   (29) } \\ & = \sum_ {p, m} \frac {(- 1) ^ {p}}{p} \psi_ {p, m} (x) \quad \text { by   (25) } \\ & = \sum_ {p} \frac {(- 1) ^ {p}}{p} \psi_ {p} (x), \end{array}
\]

whence (i). Similarly, for \(t \neq 0\) ,

\[
\begin{array}{l} \mathrm{L} (t x) = t \mathrm{L} ((t x) ^ {[ t ^ {- 1} ]}) \\ = t \mathrm{L} \left(\sum_ {m \geqslant 0} \sum_ {1 \leqslant r \leqslant m} t ^ {m - r} \phi_ {r, m} (x)\right) \quad \text { by } (24). \end{array}
\]

Equating the terms of first degree in t, we obtain

\[
x = \mathrm{L} \Bigl (\sum_ {m \geqslant 0} \phi_ {m, m} (x) \Bigr)
\]

whence

\[
\begin{array}{r l} \mathrm{E} (x) & = \sum_ {m \geqslant 0} \phi_ {m, m} (x) \\ & = \sum_ {m \geqslant 0} \frac {1}{m !} \psi_ {m, m} (x) \quad \text { by   (26) }. \end{array}
\]

PROPOSITION 4. For the chart of G used to be canonical, it is necessary and sufficient that \(\psi_j = 0\) for \(j \geqslant 2\) .

This is sufficient by Proposition 3. Suppose that the chart is canonical and that \(\psi_{i} = 0\) for \(2 \leqslant i < n\) . Then \(nx = x^{[n]} = \sum_{i=0}^{n} \binom{n}{i} \psi_{i}(x) = nx + \psi_{n}(x)\) , whence \(\psi_{n} = 0\) . Hence \(\psi_{j} = 0\) for \(j \geqslant 2\) by induction on \(j\) .

\section{REAL AND COMPLEX LIE GROUPS}

In this paragraph, K is assumed to be equal to R or C.

\subsection{PASSAGE FROM LIE ALGEBRA MORPHISMS TO LIE GROUP MORPHISMS}

Lemma 1. Let G be a simply connected† topological group, W a symmetric connected open neighbourhood of e, H a group and f a mapping of W³ into H such that

\[
f (x y z) = f (x) f (y) f (z)
\]

for \(x, y, z\) in W. There exists a morphism \(f'\) of G into H such that \(f' \mid \mathrm{W} = f \mid \mathrm{W}\) . For \((g, h) \in \mathrm{G} \times \mathrm{H}\) and U an open neighbourhood of \(e\) in W, let \(\mathrm{A}(g, h, \mathrm{U})\) be the set of \((gu, hf(u)) \in \mathrm{G} \times \mathrm{H}\) for \(u \in \mathrm{U}\) . Then \((g, h) \in \mathrm{A}(g, h, \mathrm{U})\) and \(\mathrm{A}(g, h, \mathrm{U}_1) \cap \mathrm{A}(g, h, \mathrm{U}_2) = \mathrm{A}(g, h, \mathrm{U}_1 \cap \mathrm{U}_2)\) . Let \((s, t) \in \mathrm{A}(g, h, \mathrm{U})\) ; then \(s = gu\) and \(t = hf(u)\) for \(u \in \mathrm{U}\) ; there exists an open neighbourhood \(\mathrm{U}'\) of \(e\) in W such that \(u\mathrm{U}' \in \mathrm{U}\) ; then, for \(u' \in \mathrm{U}'\) ,

\[
(s u ^ {\prime}, t f (u ^ {\prime})) = (g u u ^ {\prime}, h f (u u ^ {\prime})) \in \mathrm{A} (g, h, \mathrm{U})
\]

and hence \(\mathrm{A}(s,t,\mathrm{U}^{\prime})\subset \mathrm{A}(g,h,\mathrm{U})\) . It then follows that the \(\mathrm{A}(g,h,\mathrm{U})\) form the base of a topology on \(\mathbf{G}\times \mathbf{H}\) . We shall denote by \(\mathbf{Y}\) the set \(\mathbf{G}\times \mathbf{H}\) with this topology and denote by \(p\) the canonical projection of \(\mathbf{Y}\) onto \(\mathbf{G}\) , which is open. The restriction of \(p\) to \(\mathrm{A}(g,h,\mathrm{U})\) is a homeomorphism of \(\mathrm{A}(g,h,\mathrm{U})\) onto \(g\mathrm{U}\) . Hence \((\mathrm{Y},p)\) is a covering space of \(\mathbf{G}\) . Let \(\mathrm{Y}_0\) be the subgroup of \(\mathrm{Y}\) generated by \(\mathrm{A}(e,e,\mathrm{W})\) and let \(\mathfrak{B}\) be the set of \(\mathrm{A}(e,e,\mathrm{U})\) . Clearly \(\mathfrak{B}\) satisfies conditions \((\mathrm{GV}_{\mathrm{I}}^{\prime})\) and \((\mathrm{GV}_{\mathrm{II}}^{\prime})\) of General Topology, Chapter III, § 1, no. 2. The set \(\mathrm{Y}_0^{\prime}\) of \(y\in \mathrm{Y}_0\) such that the mappings \(z\mapsto yzy^{-1}\) and \(z\mapsto y^{-1}zy\) of \(\mathrm{Y}_0\) into \(\mathrm{Y}_0\) are continuous at \((e,e)\) is a subgroup of \(\mathrm{Y}_0\) . Let \(w\in W\) . The mapping \(w'\mapsto ww'w^{-1}\) of \(\mathrm{W}\) into \(\mathrm{G}\) is continuous and hence the mapping

\[
\left(w ^ {\prime}, f \left(w ^ {\prime}\right)\right) \mapsto \left(w w ^ {\prime} w ^ {- 1}, f \left(w w ^ {\prime} w ^ {- 1}\right)\right)
\]

of \(\mathrm{A}(e,e,\mathrm{W})\) into \(\mathrm{Y}_0\) is continuous at \((e,e)\). Now

\[
f \left(w w ^ {\prime} w ^ {- 1}\right) = f (w) f \left(w ^ {\prime}\right) f \left(w ^ {- 1}\right)
\]

\[
\left(w w ^ {\prime} w ^ {- 1}, f \left(w w ^ {\prime} w ^ {- 1}\right)\right) = (w, f (w)) \left(w ^ {\prime}, f \left(w ^ {\prime}\right)\right) \left(w, f (w)\right) ^ {- 1}
\]

As \(w^{-1} \in W\) , we see that \((w, f(w)) \in Y_0'\) . Thus, \(A(e, e, W) \subset Y_0'\) , so that \(Y_0' = Y_0\) . The group \(Y_0\) , with the filter base \(\mathfrak{B}\) , therefore satisfies condition \((GV_{\mathrm{III}}')\) of General Topology, Chapter III, § 1, no. 2. As

\[
(g, h). \mathrm{A} (e, e, \mathrm{U}) = \mathrm{A} (g, h, \mathrm{U}),
\]

\(Y_{0}\) is a topological group which is connected since \(A(e, e, W)\) is connected. Then \(p(Y_{0})\) is an open subgroup of G, whence \(p(Y_{0}) = G\) since G is connected. The kernel of \(p \mid Y_{0}\) is discrete. As G is simply connected, \(p \mid Y_{0}\) is a homeomorphism of \(Y_{0}\) onto G. Therefore \(Y_{0}\) is the graph of a morphism \(f'\) of G into H. For \(g \in W\) , \((g, f(g)) \in A(e, e, W) \subset Y_{0}\) , whence \(f(g) = f'(g)\) .

THEOREM 1. Let G and H be Lie groups and h a continuous morphism of L(G) into L(H). Suppose that G is simply connected. Then there exists one and only one Lie group morphism \(\phi\) of G into H such that \(h = \mathrm{L}(\phi)\) .

The existence of \(\phi\) follows from Lemma 1 and §4, no. 1, Theorem 1 (i). The uniqueness of \(\phi\) follows from §4, no. 1, Theorem 1 (ii) and the fact that \(G\) is connected.

COROLLARY. Let G be a finite-dimensional simply connected Lie group. There exists a finite-dimensional analytic linear representation of G whose kernel is discrete.

There exist (Chapter I, § 7, Theorem 2) a finite-dimensional vector space E and an injective morphism \(h\) of \(L(G)\) into the Lie algebra \(\operatorname{End}(E)\) . By Theorem 1, there exists a morphism \(\phi\) of G into \(\mathbf{GL}(E)\) such that \(L(\phi) = h\) . Hence \(\phi\) is an immersion and therefore its kernel is discrete.

Remarks. (1) There exist finite-dimensional simply connected Lie groups which have no finite-dimensional injective analytic linear representation (Exercise 2).

\begin{enumerate}
\def\labelenumi{(\arabic{enumi})}
\setcounter{enumi}{1}
\tightlist
\item
  There exist finite-dimensional connected Lie groups G such that every finite-dimensional analytic linear representation of G has non-discrete kernel (Exercises 3 and 4).
\end{enumerate}

\subsection{INTEGRAL SUBGROUPS}

DEFINITION 1. Let G be a Lie group. An integral subgroup of G is a subgroup H with a connected Lie group structure such that the canonical injection of H into G is an immersion.

A one-parameter subgroup of \(G\) is a 1-dimensional integral subgroup of \(G\) .

Let H be an integral subgroup of G and i the canonical injection of H into G. Then \(L(i)\) defines an isomorphism of \(L(H)\) onto a Lie subalgebra of \(L(G)\) admitting a topological supplement. \(L(H)\) is identified with its image under \(L(i)\) .

Examples. (1) A connected Lie subgroup of \(\mathbf{G}\) is an integral subgroup of \(\mathbf{G}\) .

\begin{enumerate}
\def\labelenumi{(\arabic{enumi})}
\setcounter{enumi}{1}
\item
  Suppose that \(\mathbf{G}\) is finite-dimensional. Let \(\mathbf{H}\) be a subgroup of \(\mathbf{G}\) ; we give it the structure induced by the Lie group structure on \(\mathbf{G}\) (§ 4, no. 5, Definition 3). Then its identity component \(\mathrm{H}_0\) is an integral subgroup of \(\mathbf{G}\) and the subalgebra tangent to \(\mathbf{H}\) at \(e\) is \(\mathrm{L}(\mathrm{H}_0)\) (§ 4, no. 5, Proposition 9 (ii)).
\item
  Let \(\mathbf{G}\) be a complex Lie group, \(\mathbf{H}\) an integral subgroup of \(\mathbf{G}\) and \(\mathbf{G}_1\) (resp. \(\mathbf{H}_1\) ) the underlying real Lie group of \(\mathbf{G}\) (resp. \(\mathbf{H}\) ). Then \(\mathbf{H}_1\) is an integral subgroup of \(\mathbf{G}_1\) and \(\mathbf{L}(\mathbf{H}_1)\) is the underlying real Lie algebra of \(\mathbf{L}(\mathbf{H})\) .
\end{enumerate}

THEOREM 2. Let G be a Lie group.

\begin{enumerate}
\def\labelenumi{(\roman{enumi})}
\item
  The mapping \(\mathbf{H} \mapsto \mathbf{L}(\mathbf{H})\) is a bijection of the set of integral subgroups of \(\mathbf{G}\) onto the set of Lie subalgebras of \(\mathbf{L}(\mathbf{G})\) admitting a topological supplement.
\item
  Let \(\mathbf{H}\) be an integral subgroup of \(\mathbf{G}\) . Every connected Lie subgroup germ of \(\mathbf{G}\) with Lie algebra \(\mathbf{L}(\mathbf{H})\) is an open submanifold of \(\mathbf{H}\) which generates \(\mathbf{H}\) .
\end{enumerate}

\begin{enumerate}
\def\labelenumi{(\alph{enumi})}
\item
  Let \(\mathfrak{b}\) be a Lie subalgebra of \(\mathrm{L}(\mathrm{G})\) admitting a topological supplement. Let \(\mathrm{H}_{1}\) be a Lie subgroup germ of \(\mathrm{G}\) such that \(\mathrm{L}(\mathrm{H}_{1}) = \mathfrak{b}\) (§ 4, Theorem 3). \(\mathrm{H}_{1}\) can be chosen so that it is connected. Let \(\mathrm{H}\) be the subgroup of \(\mathrm{G}\) generated by \(\mathrm{H}_{1}\) . There exists (§ 1, Corollary to Proposition 22) a Lie group structure on \(\mathrm{H}\) such that \(\mathrm{H}_{1}\) is an open submanifold of \(\mathrm{H}\) and the canonical injection of \(\mathrm{H}\) into \(\mathrm{G}\) is an immersion. As \(\mathrm{H}_{1}\) is connected, \(\mathrm{H}\) is connected and is hence an integral subgroup of \(\mathrm{G}\) . Then \(\mathrm{L}(\mathrm{H}) = \mathrm{L}(\mathrm{H}_{1}) = \mathfrak{b}\) . This proves that the mapping considered in (i) is surjective.
\item
  Let H be an integral subgroup of G and \(N_{1}\) a connected Lie subgroup germ of G with Lie algebra L(H). As the canonical injection of H into G is an immersion, there exists an open subgroup germ \(H_{1}\) of H which is at the same time a submanifold of G and hence a Lie subgroup germ of G with Lie algebra L(H). On the other hand, let N be the subgroup of G generated by \(N_{1}\) ; by part (a) of the proof, it has the structure of an integral subgroup of G such that \(N_{1}\) is an open submanifold of N. By §4, Theorem 3, \(H_{1} \cap N_{1}\) is open in \(H_{1}\) and \(N_{1}\) . Hence the subgroup of G generated by \(H_{1} \cap N_{1}\) is equal on the one hand to H and on the other to N. Therefore the Lie groups H and N are equal. This proves (ii) and also proves that the mapping considered in (i) is injective.
\end{enumerate}

Remark 1. Let H be an integral subgroup of G. Let Y be the left foliation of G associated with L(H). If \(g \in G\) , let gH be given the manifold structure derived from that on H by \(\gamma(g)\) . By Differentiable and Analytic Manifolds, R, 9.3.2, the canonical injection of gH into Y is a morphism. This morphism is étale. Hence the maximal connected leaves of Y are the left cosets modulo H.

PROPOSITION 1. Let G and M be Lie groups, H an integral subgroup of G and \(\phi\) a morphism of M into G such that \(\mathrm{L}(\phi)(\mathrm{L}(\mathrm{M})) \subset \mathrm{L}(\mathrm{H})\) . Suppose that M is connected. Then \(\phi(\mathrm{M}) \subset \mathrm{H}\) and \(\phi\) , considered as a mapping of M into H, is a Lie group morphism.

In the notation of Remark 1, \(\phi\) is a morphism of M into Y (Differentiable and Analytic Manifolds, R, 9.3.2) and hence \(\phi(M) \subset H\) since M is connected.

COROLLARY 1. Let G and H be Lie groups, \(\phi\) a Lie group morphism of G into H, N the kernel of \(\phi\) and \(h = L(\phi)\) . Suppose that G is connected and that H is finite-dimensional.

\begin{enumerate}
\def\labelenumi{(\roman{enumi})}
\item
  N is a Lie subgroup of G and L(N) = Ker h.
\item
  Let \(\mathbf{H}'\) be the integral subgroup of \(\mathbf{H}\) with Lie algebra \(\operatorname{Im} h\) . Then \(\phi(\mathbf{G}) = \mathbf{H}'\) .
\item
  The mapping of G/N into H' derived from \(\phi\) when passing to the quotient is a Lie group isomorphism.
\end{enumerate}

\((i)\) has already been proved (§ 3, no. 8, Proposition 28).

Let \(\psi\) be the Lie group morphism of G/N into H derived from \(\phi\) when passing to the quotient; it is an immersion (§3, no. 8, Proposition 28). By Proposition 1, \(\psi\) is a Lie group morphism of G/N into H'. This morphism is étale and hence \(\psi\) (G/N) = H' since H' is connected; this proves (ii). Then \(\psi\colon G/N \to H'\) is bijective and is a Lie group isomorphism, which proves (iii).

COROLLARY 2. Let \(\mathbf{G}\) be a Lie group and \(\mathrm{H}_1\) and \(\mathrm{H}_2\) integral subgroups of \(\mathbf{G}\) . If \(\mathrm{L}(\mathrm{H}_2) \subset \mathrm{L}(\mathrm{H}_1)\) , then \(\mathrm{H}_2\) is an integral subgroup of \(\mathrm{H}_1\) .

Let \(i_1: \mathrm{H}_1 \to \mathrm{G}\) , \(i_2: \mathrm{H}_2 \to \mathrm{G}\) be the canonical injections. Then

\[
\mathrm{L} (i _ {2}) (\mathrm{L} (\mathrm{H} _ {2})) = \mathrm{L} (\mathrm{H} _ {2}) \subset \mathrm{L} (\mathrm{H} _ {1}).
\]

By Proposition 1, \(i_{2}\) is an analytic mapping of \(H_{2}\) into \(H_{1}\) and even an immersion of \(H_{2}\) into \(H_{1}\) since \(\mathrm{L}(i_{2})\) is an isomorphism of \(\mathrm{L}(\mathrm{H}_{2})\) onto a subalgebra of \(\mathrm{L}(\mathrm{H}_{1})\) admitting a topological supplement.

COROLLARY 3. Let \(\mathbf{G}\) be a finite-dimensional Lie group and \((\mathrm{H}_i)_{i\in I}\) a family of Lie subgroups of \(\mathbf{G}\) . Then \(\mathrm{H} = \bigcap_{i\in I}\mathrm{H}_i\) is a Lie subgroup of \(\mathbf{G}\) and

\[
\mathrm{L} (\mathrm{H}) = \bigcap_ {i \in \mathrm{I}} \mathrm{L} (\mathrm{H} _ {i}).
\]

There exists a finite subset J of I such that \(\bigcap_{i\in J}L(H_i)\) is equal to the intersection M of all the \(\mathrm{L}(\mathrm{H}_i)\) . We know that \(\mathrm{H}^* = \bigcap_{i\in J}\mathrm{H}_i\) is a Lie subgroup such that \(\mathrm{L}(\mathrm{H}^*) = \mathrm{M}\) (§ 3, no. 8, Corollary 2 to Proposition 29). Let \(\mathrm{H}_0\) be the identity component of \(\mathrm{H}^*\) . It is a Lie subgroup of G and \(\mathrm{L}(\mathrm{H}_0) = \mathrm{M}\) . By Corollary 2, \(\mathrm{H}_0\subset \mathrm{H}_i\) for all \(i\) and hence \(\mathrm{H}_0\subset \mathrm{H}\subset \mathrm{H}^*\) , whence the corollary.

COROLLARY 4. Let G be a finite-dimensional connected Lie group. The following conditions are equivalent:

\begin{enumerate}
\def\labelenumi{(\roman{enumi})}
\item
  G is unimodular (Integration, Chapter VII, § 1, no. 3, Definition 3);
\item
  det Ad g = 1 for all g ∈ G;
\item
  \(\operatorname{Tr} \operatorname{ad} a = 0\) for all \(a \in \mathrm{L}(\mathrm{G})\) .
\end{enumerate}

The mapping \(g \mapsto \operatorname{det} \operatorname{Ad} g\) is a morphism \(\phi\) of \(G\) into \(K^*\) . By § 3, Propositions 35 (no. 10) and 44 (no. 12), \(L(\phi)a = Tr\operatorname{ad} a\) for all \(a \in L(G)\) . Clearly \(\operatorname{Im} L(\phi) = \{0\}\) or \(K\) . In the first (resp. second) case, \(\operatorname{Im} \phi = \{1\}\) (resp. \(\operatorname{Im} \phi = K^*\) ) by Corollary 1 and hence \(G\) is unimodular (resp. not unimodular) by § 3, no. 16, Corollary to Proposition 55.

PROPOSITION 2. Let \(\mathbf{G}\) be a finite-dimensional Lie group and \(\mathbf{H}\) an integral subgroup of \(\mathbf{G}\) . The following conditions are equivalent:

\begin{enumerate}
\def\labelenumi{(\roman{enumi})}
\item
  H is closed;
\item
  the topology on H is induced by that on G;
\item
  H is a Lie subgroup of G.
\end{enumerate}

\((i) \Rightarrow (iii)\): this follows from § 1, Propositions 2 (iv) (no. 1) and 14 (iii) (no. 7).

\((iii) \Rightarrow (ii)\): obvious.

\((ii) \Rightarrow (i)\): if the topology on H is induced by that on G, H is closed because H is complete (§ 1, no. 1, Proposition 1).

PROPOSITION 3. Let G be a Lie group, H an integral subgroup of G, M a non-empty connected analytic manifold, f a mapping of M into G and \(r \in N_{K}\) . Consider the following conditions:

\begin{enumerate}
\def\labelenumi{(\roman{enumi})}
\item
  \(f\) is of class \(\mathbf{C}^r\) and \(f(\mathbf{M})\subset \mathbf{H}\) ;
\item
  \(f(\mathbf{M})\subset \mathbf{H}\) and \(f\) , considered as a mapping of \(\mathbf{M}\) into \(\mathbf{H}\) , is of class \(\mathbf{C}^r\)
\item
  \(f\) is of class \(\mathbf{C}^r\) , \(f(\mathbf{M})\) meets \(\mathbf{H}\) and the image of \(\mathrm{T}_m(\mathbf{M})\) is contained in \(f(m)\) . \(\mathbf{L}(\mathbf{H})\) for all \(m \in \mathbf{M}\) .
\end{enumerate}

Then (ii) \(\Leftrightarrow\) (iii) \(\Rightarrow\) (i). If the topology on H admits a countable base, the three conditions are equivalent.

\((ii) \Rightarrow (i)\) and \((ii) \Rightarrow (iii)\): obvious.

\((iii) \Rightarrow (ii)\): suppose that condition (iii) holds. By Differentiable and Analytic Manifolds, R, 9.2.8, \(f\) is a morphism of class \(\mathbf{C}^r\) of M into the left foliation associated with \(\mathbf{L}(\mathbf{H})\) . As M is connected, \(f(\mathbf{M}) \subset \mathbf{H}\) .

If the topology on H admits a countable base, condition (i) implies that f is a mapping of class \(C^{r}\) of M into H (Differentiable and Analytic Manifolds, R, 9.2.8); hence (i) \(\Rightarrow\) (ii).

COROLLARY 1. Let G be a finite-dimensional Lie group and H an integral subgroup of G. Then the Lie subalgebra tangent to H at e (§ 4, no. 5, Definitions 2 and 3) is L(H) and the Lie group structure on H is the structure induced by that on G.

As H is connected and finite-dimensional, its topology admits a countable base.

COROLLARY 2. Let G be a Lie group and \(H_{1}\) and \(H_{2}\) integral subgroups of G. Suppose that the topology on \(H_{1}\) admits a countable base. Then

\[
\mathrm{H} _ {2} \subset \mathrm{H} _ {1} \Leftrightarrow \mathrm{L} (\mathrm{H} _ {2}) \subset \mathrm{L} (\mathrm{H} _ {1})
\]

and, if these conditions hold, \(\mathbf{H}_2\) is an integral subgroup of \(\mathbf{H}_1\) .

The last assertion and the implication \(\mathrm{L}(\mathrm{H}_{2})\subset\mathrm{L}(\mathrm{H}_{1})\Rightarrow\mathrm{H}_{2}\subset\mathrm{H}_{1}\) follow from Corollary 2 to Proposition 1. The converse implication follows from Proposition 3.

COROLLARY 3. Let G be a Lie group and \(H_{1}\) and \(H_{2}\) integral subgroups of G whose topology admits a countable base. If \(H_{1}\) and \(H_{2}\) have the same underlying set, the Lie group structures on \(H_{1}\) and \(H_{2}\) are the same.

This follows from Corollary 2.

Remark 2. Let G be a finite-dimensional Lie group. Let H be a subgroup of G. We shall say, by an abuse of language, that H is an integral subgroup of G if there exists a Lie group structure S on H such that H, together with S, is an integral subgroup of G. By Corollary 3 to Proposition 3, if S exists, S is unique.

Remark 3. Let V be a manifold of class \(\mathbf{C}^r\) . Let M be a subset of V and \(x\) and \(y\) elements of M. Consider the following property:

\(P_{M,x,y}\) : there exist I, \(x_{0}, x_{1}, \ldots, x_{n}, f_{1}, \ldots, f_{n}\) such that: (a) I is a connected open subset of K; (b) \(x_{0}, \ldots, x_{n}\) are in M, \(x_{0} = x, x_{n} = y\) ; (c) for \(1 \leqslant i \leqslant n\) , \(f_{i}\) is a mapping of class \(C^{r}\) of I into V which takes the values \(x_{i-1}\) and \(x_{i}\) and \(f_{i}(I) \subset M\) .

We shall say that M is a \(C^{r}\) -connected subset of V if, for all elements x, y of M, the property \(P_{M,x,y}\) holds.

PROPOSITION 4. Let G be a finite-dimensional Lie group and H a subgroup of G. Let \(r \in \mathbf{N}_{\mathbb{K}}\) . The following conditions are equivalent:

\begin{enumerate}
\def\labelenumi{(\roman{enumi})}
\item
  H is an integral subgroup of G;
\item
  with the Lie group structure induced by that on G, H is connected;
\item
  H is \(\mathbf{C}^r\) -connected.
\end{enumerate}

\((ii) \Rightarrow (i)\): obvious.

\((i) \Rightarrow (iii)\): suppose that H has a Lie group structure such that H is an integral subgroup of G. Using the notation of Remark 3, the set of \(y \in H\) such that property \(\mathrm{P}_{\mathrm{H},e,y}\) holds is an open subgroup of H. As H is connected, this subgroup is equal to H and hence condition (iii) is satisfied.

\((iii) \Rightarrow (ii)\): suppose that condition (iii) is satisfied and let H be given the structure induced by the Lie group structure on G. Let \(\mathfrak{h}\) be the subalgebra tangent to H at \(e\) . The identity component \(\mathrm{H}_0\) of H is an integral subgroup of G such that \(\mathrm{L}(\mathrm{H}_0) = \mathfrak{h}\) . We show that \(\mathrm{H} = \mathrm{H}_0\) . It suffices to prove the following: let I be a connected open subset of K, \(f\) a mapping of class \(\mathbf{C}^r\) of I into G such that \(f(\mathrm{I}) \subset \mathrm{H}\) and \(\lambda\) and \(\mu\) two points of I; if \(f(\lambda) \in \mathrm{H}_0\) , then \(f(\mu) \in \mathrm{H}_0\) . But, for all \(\nu \in \mathrm{I}\) , \((\mathrm{T}_{\nu}f)(\mathrm{K}) \subset f(\nu)\mathfrak{h}\) by definition of \(\mathfrak{h}\) , so that our assertion follows from Proposition 3.

Remark 4. If \(\mathbf{K} = \mathbf{R}\) , the integral subgroups of \(\mathbf{G}\) can also be characterized as the subgroups which, with the topology induced by that on \(\mathbf{G}\) , are arcwise connected (§ 8, Exercise 4). However, subgroups may be connected but not integral (Commutative Algebra, Chapter VI, § 9, Exercise 2).

COROLLARY. Let G be a finite-dimensional Lie group and \(H_{1}\) and \(H_{2}\) two integral subgroups of G. The subgroup of G generated by \(H_{1}\) and \(H_{2}\) and the subgroup ( \(H_{1}, H_{2}\) ) of G are integral subgroups of G.

The subgroup (G, G) of G is not always closed (§ 9, Exercise 6).

Recall (§ 3, no. 11, Corollary 5 to Proposition 41) that if a is a finite-dimensional algebra, \(\operatorname{Aut}(a)\) is a Lie subgroup of \(\mathbf{GL}(a)\) and that \(\mathrm{L}(\mathrm{Aut}(a))\) is the Lie algebra of derivations of \(a\) .

DEFINITION 2. Let \(\mathfrak{a}\) be a finite-dimensional Lie algebra. Let \(\operatorname{Ad}(\mathfrak{a})\) or \(\operatorname{Int}(\mathfrak{a})\) denote the integral subgroup of \(\operatorname{Aut}(\mathfrak{a})\) with Lie algebra \(\operatorname{ad}(\mathfrak{a})\) . The elements of this group are called inner automorphisms of \(\mathfrak{a}\) .

By transport of structure, \(\operatorname{ad}(a)\) is invariant under \(\operatorname{Aut}(a)\) and hence \(\operatorname{Int}(a)\) is normal in \(\operatorname{Aut}(a)\) . By § 4, no. 4, Corollary 1 to Proposition 8 and the fact that \(\operatorname{Int}(a)\) is connected, the elements of \(\operatorname{Int}(a)\) are the finite products of automorphisms of the form \(\exp ad x\) where \(x \in a\) . In general, \(\operatorname{Int}(a)\) is not a Lie subgroup of \(\operatorname{Aut}(a)\) (Exercise 14).

\subsection{PASSAGE FROM LIE ALGEBRAS TO LIE GROUPS}

THEOREM 3. (i) If L is a finite-dimensional Lie algebra, there exists a simply connected Lie group G such that L(G) is isomorphic to L.

\begin{enumerate}
\def\labelenumi{(\roman{enumi})}
\setcounter{enumi}{1}
\tightlist
\item
  Let \(G_{1}\) and \(G_{2}\) be two connected Lie groups, with \(G_{1}\) simply connected. Let f be an isomorphism of \(\mathrm{L}(\mathrm{G}_{1})\) onto \(\mathrm{L}(\mathrm{G}_{2})\) , \(\phi\) the morphism of \(G_{1}\) into \(G_{2}\) such that \(\mathrm{L}(\phi)=f\) and N the kernel of \(\phi\) . Then N is a discrete subgroup of the centre of \(G_{1}\) and the morphism of \(G_{1}/N\) into \(G_{2}\) derived from \(\phi\) is a Lie group isomorphism. If \(G_{2}\) is simply connected, \(\phi\) is an isomorphism.
\end{enumerate}

Let \(\mathbf{L}\) be a finite-dimensional Lie algebra. There exists a finite-dimensional vector space \(\mathbf{E}\) such that \(\mathbf{L}\) can be identified with a Lie subalgebra of \(\operatorname{End}(\mathbf{E})\) (Chapter I, §7, Theorem 2). Let \(\mathbf{H}\) be the integral subgroup of \(\mathbf{GL}(\mathbf{E})\) with

Lie algebra L. Let \(\hat{\mathbf{H}}\) be its universal covering (§ 1, no. 9, Remark). Then \(\mathbf{L}(\hat{\mathbf{H}})\) is isomorphic to L, whence (i).

Let \(G_{1}, G_{2}, \tilde{f}, \phi, N\) be as in (ii). Then \(\phi\) is étale, hence \(\phi(G_{1})\) is an open subgroup of \(G_{2}\) and hence \(\phi(G_{1}) = G_{2}\) . On the other hand, N is discrete and hence contained in the centre of \(G_{1}\) (Integration, Chapter VII, § 3, Lemma 4). Clearly the morphism of \(G_{1}/N\) onto \(G_{2}\) derived from \(\phi\) is a Lie group isomorphism. If \(G_{2}\) is simply connected, every étale mapping of \(G_{1}\) onto \(G_{2}\) is injective and hence \(N = \{e\}\) .

PROPOSITION 5. Let G be a connected real Lie group. Suppose that L(G) has a complex normable Lie algebra structure L' compatible with its real normable Lie algebra structure. There exists on G one and only one complex Lie group structure compatible with the real Lie group structure and with Lie algebra L'.

By § 4, no. 2, Corollary 2 to Theorem 2, it suffices to prove that the structure of \(L'\) is invariant under Ad G. Let \(\phi\) be an exponential mapping of G. By § 4, no. 4, Corollary 3 (i) to Proposition 8, there exists a neighbourhood V of 0 in \(L(G)\) such that the structure of \(L'\) is invariant under \(Ad\phi(V)\) . But \(Ad\phi(V)\) generates G because G is connected.

The conclusion of Proposition 5 is not necessarily true if G is not assumed to be connected (Exercise 7).

PROPOSITION 6. Let G be a connected complex Lie group. If G is compact, G is commutative.

The holomorphic mapping \(g \mapsto \operatorname{Ad} g\) of \(\mathbf{G}\) into \(\mathcal{L}(\mathbf{L}(\mathbf{G}))\) is constant (Differentiable and Analytic Manifolds, R, 3.3.7) and hence ad \(a = 0\) for all \(a \in \mathbf{L}(\mathbf{G})\) (§ 3, no. 12, Proposition 44). Hence \(\mathbf{G}\) is commutative (§ 4, Corollary 3 to Theorem 1).

\subsection{EXPONENTIAL MAPPING}

THEOREM 4. Let \(\mathbf{G}\) be a Lie group. There exists one and only one exponential mapping of \(\mathbf{G}\) defined on \(\mathrm{L}(\mathbf{G})\) .

There exist a convex open neighbourhood \(\mathbf{U}\) of 0 in \(\mathrm{L}(\mathbf{G})\) and an exponential mapping of \(\mathbf{G}\) defined on \(\mathbf{U}\) . It can be assumed, by choosing \(\mathbf{U}\) sufficiently small, that

\[
\phi ((\lambda + \lambda^ {\prime}) a) = \phi (\lambda a) \phi (\lambda^ {\prime} a)\tag{1}
\]

for \(a \in \mathrm{L}(\mathbf{G})\) , \(\lambda, \lambda'\) in \(\mathbf{K}\) , \(\lambda a, \lambda'a\) , \((\lambda + \lambda')a\) in \(\mathbf{U}\) .

Let \(a \in \mathrm{L}(\mathrm{G})\) . There exists an integer \(n > 0\) such that \(\frac{1}{n} a \in \mathrm{U}\) . If \(m\) is another integer such that \(\frac{1}{m} a \in \mathrm{U}\) , then \(\frac{1}{nm} a \in \mathrm{U}\) and relation (1) implies

\[
\phi \left(\frac {1}{n} a\right) = \left(\phi \left(\frac {1}{n m} a\right)\right) ^ {m}, \quad \phi \left(\frac {1}{m} a\right) = \left(\phi \left(\frac {1}{n m} a\right)\right) ^ {n};
\]

hence \(\left(\phi\left(\frac{1}{n}a\right)\right)^n = \left(\phi\left(\frac{1}{m}a\right)\right)^m\) . There exists an extension \(\psi: \mathrm{L}(G) \to G\) of \(\phi\) such that \(\psi(a) = \left(\phi\left(\frac{1}{n}a\right)\right)^n\) for \(a \in \mathrm{L}(G)\) and \(n\) an integer \(>0\) such that \(\frac{1}{n}a \in \mathrm{U}\) . Clearly \(\psi\) is analytic and is an exponential mapping of \(G\) . If \(\psi': \mathrm{L}(G) \to G\) is an exponential mapping of \(G\) , \(\psi\) and \(\psi'\) coincide on a neighbourhood of 0 and hence are equal since \(\mathrm{L}(G)\) is connected.

Henceforth when we speak of the exponential mapping of G, we shall mean the mapping considered in Theorem 4. It will be denoted by \(\exp_{G}\) or exp if there is no risk of confusion.

Example. Let A be a complete normed unital associative algebra. Then \(\exp_{\mathbf{A}}\) is the exponential mapping defined in Chapter II, § 7, no. 3.

PROPOSITION 7. Let G be a Lie group and a an element of L(G). The mapping \(\lambda \mapsto \exp(\lambda a)\) of K into G is the unique morphism \(\phi\) of the Lie group K into G such that \((\mathrm{T}_0\phi)1 = a\) .

The mappings \((\lambda, \lambda') \mapsto \exp(\lambda a) \exp(\lambda' a)\) and \((\lambda, \lambda') \mapsto \exp(\lambda + \lambda') a\) of \(K \times K\) into \(G\) are analytic and coincide on a neighbourhood of \((0, 0)\) . As \(K \times K\) is connected, these mappings are equal. Hence \(\phi: \lambda \mapsto \exp(\lambda a)\) is a Lie group morphism of \(K\) into \(G\) . The tangent mapping at \(0\) to \(\lambda \mapsto \lambda a\) is the mapping \(\lambda \mapsto \lambda a\) ; and \(T_0(\exp) = \mathrm{Id}_{\mathrm{L}(\mathrm{G})}\) ; hence \((T_0\phi)1 = a\) . The uniqueness assertion of the proposition follows from Theorem 1.

PROPOSITION 8. Let G be a Lie group. For all x, y in L(G) and n an integer,

\begin{enumerate}
\def\labelenumi{(\arabic{enumi})}
\setcounter{enumi}{1}
\tightlist
\item
\end{enumerate}

\[
\exp (x + y) = \lim _ {n \rightarrow + \infty} \left(\left(\exp \frac {1}{n} x\right)\left(\exp \frac {1}{n} y\right)\right) ^ {n}\tag{2}
\]

\[
\exp [ x, y ] = \lim _ {n \rightarrow + \infty} \left(\left(\exp \frac {1}{n} x\right)\left(\exp \frac {1}{n} y\right)\left(\exp \frac {1}{n} x\right) ^ {- 1} \left(\exp \frac {1}{n} y\right) ^ {- 1}\right) ^ {n ^ {2}}\tag{3}.
\]

By Proposition 7, this follows from Proposition 4 of § 4, no. 3, taking \(\lambda = \frac{1}{n}\) .

PROPOSITION 9. Let G be a complex Lie group and G' the underlying real Lie group. Then \(\exp_{G} = \exp_{G'}\) .

This follows from Proposition 5 of §4, no. 3 and the analyticity of \(\exp_{G}\) and \(\exp_{G'}\) .

PROPOSITION 10. Let G and H be Lie groups and \(\phi\) a morphism of G into H.

\begin{enumerate}
\def\labelenumi{(\roman{enumi})}
\item
  \(\phi \circ \exp_{\mathrm{G}} = \exp_{\mathrm{H}} \circ \mathrm{L}(\phi)\) .
\item
  If G is an integral subgroup of H, then \(\exp_{G} = \exp_{H} | L(G)\) .
\end{enumerate}

The two sides of the equality (i) are analytic mappings of L(G) into H which coincide in a neighbourhood of 0 (§ 4, Proposition 8, no. 4) and hence are equal. Assertion (ii) is a special case of (i).

COROLLARY 1. Let G be a Lie group, G' a Lie subgroup of G and \(a \in \mathrm{L}(\mathrm{G})\) . The following conditions are equivalent:

\begin{enumerate}
\def\labelenumi{(\roman{enumi})}
\item
  \(a \in \mathbf{L}(\mathbf{G}')\) ;
\item
  \(\exp (\lambda a)\in \mathbf{G}'\) for all \(\lambda \in \mathbf{K}\) and \(|\lambda |\) sufficiently small;
\item
  \(\exp(\lambda a) \in \mathbf{G}'\) for all \(\lambda \in K\) .
\end{enumerate}

The argument is as in § 4, no. 4, Corollary 1 to Proposition 8.

COROLLARY 2. Let \(\mathbf{G}\) be a Lie group, \(\mathbf{H}\) an integral subgroup of \(\mathbf{G}\) and \(a \in \mathbf{L}(\mathbf{G})\) . Consider the following conditions:

\begin{enumerate}
\def\labelenumi{(\roman{enumi})}
\item
  \(a \in \mathrm{L}(\mathrm{H})\) ;
\item
  \(\exp_{\mathbf{G}}(\lambda a)\in \mathbf{H}\) for all \(\lambda \in \mathbf{K}\) .
\end{enumerate}

Then (i) \(\Rightarrow\) (ii). If the topology of H admits a countable base, then (i) \(\Leftrightarrow\) (ii).

Let \(i\) be the canonical injection of \(\mathbf{H}\) into \(\mathbf{G}\) . If \(a \in \mathbf{L}(\mathbf{H})\) , then

\[
\exp_ {\mathrm{G}} (\lambda a) = \left(\exp_ {\mathrm{G}} \circ \mathrm{L} (i)\right) (\lambda a) = (i \circ \exp_ {\mathrm{H}}) (\lambda a) \in \mathrm{H}.
\]

Hence (i) \(\Rightarrow\) (ii). The converse when H has a countable base follows from Proposition 3.

COROLLARY 3. Let G be a Lie group, \(\rho\) an analytic linear representation of G, \(x \in L(G)\) and \(g \in G\) .

\begin{enumerate}
\def\labelenumi{(\roman{enumi})}
\item
  \(\rho(\exp x) = \exp L(\rho)x;\)
\item
  \(\mathrm{Ad}(\exp x) = \exp \mathrm{ad} x;\)
\item
  \(g(\exp x)g^{-1} = \exp(\operatorname{Ad} g.x).\)
\end{enumerate}

The argument is as in § 4, no. 4, Corollaries 2 and 3 to Proposition 8.

COROLLARY 4. Let G be a finite-dimensional connected Lie group.

\begin{enumerate}
\def\labelenumi{(\roman{enumi})}
\item
  \(\text{Int}(L(G)) = \text{Ad}(G).\)
\item
  Let \(\mathbf{Z}\) be the centre of \(\mathbf{G}\) . Then \(\mathbf{Z}\) is a Lie subgroup of \(\mathbf{G}\) whose Lie algebra is the centre of \(\mathbf{L}(\mathbf{G})\) . The mapping of \(\mathbf{G} / \mathbf{Z}\) into \(\operatorname{Int} \mathbf{L}(\mathbf{G})\) derived from \(g \mapsto \operatorname{Ad} g\) when passing to the quotient is a Lie group isomorphism.
\end{enumerate}

Assertion (i) follows from Corollary 3 (ii) and the remarks after Definition 2. Let \(g \in G\) . In order that \(\operatorname{Ad} g = \operatorname{Id}_{\operatorname{L(G)}}\) , it is necessary and sufficient that \(\operatorname{Int} g\) coincide with \(\operatorname{Id}_G\) on a neighbourhood of \(e\) (§ 4, no. 1, Theorem 1 (ii)) and hence on the whole of \(G\) ; in other words, it is necessary and sufficient that \(g \in Z\) . Then (ii) follows from Corollary 1 to Proposition 1.

DEFINITION 3. Let G be a finite-dimensional connected Lie group. The Lie group \(\operatorname{Int}(\mathbf{L}(\mathbf{G})) = \operatorname{Ad}(\mathbf{G})\) is called the adjoint group of G.

PROPOSITION 11. Let \(\mathbf{G}\) be a connected commutative Lie group.

\begin{enumerate}
\def\labelenumi{(\roman{enumi})}
\item
  exp is an étale morphism of the additive Lie group L(G) onto G.
\item
  If \(\mathbf{K} = \mathbf{R}\) and \(\mathbf{G}\) is finite-dimensional, \(\mathbf{G}\) is isomorphic to a Lie group of the form \(\mathbf{R}^p \times \mathbf{T}^q\) ( \(p, q\) integers \(\geqslant 0\) ).
\end{enumerate}

By the Hausdorff formula, \((\exp x)(\exp y) = \exp (x + y)\) for \(x,y\) sufficiently close to 0 and hence for all \(x,y\) in \(\mathrm{L}(G)\) by analytic continuation. Hence exp is a group homomorphism and is étale since

\[
\mathrm{T} _ {e} (\exp) = \mathrm{Id} _ {\mathrm{L(G)}}.
\]

Hence (i). Assertion (ii) follows from (i) and General Topology, Chapter VII, § 1, Proposition 9.

PROPOSITION 12. Let G be a Lie group and L = L(G). For all \(x \in L\) , let \(\mathrm{T}_{x}(\mathrm{L})\) be identified with L, so that the right differential \(\varpi(x)\) of exp at x is a linear mapping of L into L. For all \(x \in L\) ,

\[
\varpi (x) = \sum_ {n \geqslant 0} \frac {1}{(n + 1) !} (\mathrm{ad} x) ^ {n}.
\]

The two sides are analytic functions of \(x\) and are equal for \(x\) sufficiently close to 0 (§ 4, no. 3, Proposition 6).

Remark. \(\varpi(x) . (\mathrm{ad} x) = \exp \mathrm{ad} x - 1\) . We write, by an abuse of notation,

\[
\varpi (x) = \frac {\exp \operatorname{ad} x - 1}{\operatorname{ad} x}.
\]

COROLLARY. Let \(\mathbf{G}\) be a complex Lie group and \(x \in \mathrm{L}(\mathbf{G})\) . The tangent mapping at \(x\) to \(\exp_{\mathbf{G}}\) has kernel \(\bigoplus_{n \in \mathbf{Z}\setminus\{0\}} \operatorname{Ker}(\operatorname{ad} x - 2i\pi n)\) .

The integral function \(z \mapsto \sum_{n \geqslant 0} \frac{1}{(n + 1)!} z^n\) , equal to \(\frac{e^z - 1}{z}\) for \(z \neq 0\) , admits

as zeros the points of \(2\pi i\mathbf{Z}\setminus\{0\}\) , which are all simple zeros. The corollary then follows from Proposition 12 and the following lemma:

Lemma 2. Let E be a complex Banach space, u an element of \(\mathcal{L}(\mathrm{E})\) , S the spectrum of \(u\) in \(\mathcal{L}(\mathrm{E})\) (Spectral Theories, Chapter I, §1, no. 2) and \(f\) a holomorphic complex function on an open neighbourhood \(\Omega\) of S. Suppose that \(f\) admits in \(\Omega\) only a finite number of distinct zeros \(z_{1},\ldots ,z_{n}\) , of multiplicities \(h_1,\dots,h_n\) . Then \(\operatorname {Ker}f(u)\) is the direct sum of the \(\operatorname {Ker}(u - z_i)^{h_i}\) for \(1\leqslant i\leqslant n\) .

(For the definition of \(f(u)\) , see Spectral Theories, Chapter I, § 4, no. 8.)

There exists a holomorphic function \(g\) on \(\Omega\) , everywhere non-zero, such that \(f(z) = (z - z_1)^{h_1} \ldots (z - z_n)^{h_n} g(z)\) . Then \(g(u)g^{-1}(u) = g^{-1}(u)g(u) = 1\) and hence \(\operatorname{Ker}f(u) = \operatorname{Ker}\prod_{i=1}^{n}(u - z_i)^{h_i}\) . Considering \(\operatorname{Ker}f(u)\) as a \(\mathbf{C}[\mathbf{X}]\) -module by means of the external law \((h,x)\mapsto h(u)x\) for \(h\in \mathbf{C}[\mathbf{X}]\) , \(x\in \operatorname{Ker}f(u)\) , we see that \(\operatorname{Ker} f(u)\) is the direct sum of the \(\operatorname{Ker}(u - z_i)^{h_i}\) , using Algebra, Chapter VII, § 2, no. 1, Proposition 1.

\subsection{APPLICATION TO LINEAR REPRESENTATIONS}

PROPOSITION 13. Let G be a connected Lie group and \(\rho\) an analytic linear representation of G on a complete normable space E. Let \(\mathrm{E}_1, \mathrm{E}_2\) be two closed vector subspaces of E such that \(\mathrm{E}_2 \subset \mathrm{E}_1\) . The following conditions are equivalent:

\begin{enumerate}
\def\labelenumi{(\roman{enumi})}
\item
  \(\rho(g)x \equiv x \pmod{\mathbf{E}_2}\) for all \(g \in \mathbf{G}\) and all \(x \in \mathbf{E}_1\) ;
\item
  \(\mathrm{L}(\rho)(\mathrm{L}(\mathrm{G}))\) maps \(E_{1}\) into \(E_{2}\) .
\end{enumerate}

\[
\begin{array}{l l l} \rho (g) x \equiv x & (\mathrm{mod} \mathrm{E} _ {2}) & \text {for all \(g\in G\) and all \(x\in E_{1}\)} \\ \Leftrightarrow \rho (\exp a) x \equiv x & (\mathrm{mod} \mathrm{E} _ {2}) & \text {for all \(a\in L(G)\) and all \(x\in E_{1}\)} \\ \Leftrightarrow (\exp \mathrm{L} (\rho) a) x \equiv x & (\mathrm{mod} \mathrm{E} _ {2}) & \text {for all \(a\in L(G)\) and all \(x\in E_{1}\)}. \end{array}
\]

On the other hand, if \(u \in \mathcal{L}(\mathrm{E})\) , then

\[
\begin{array}{r l} & \exp (\lambda u) x \equiv x \pmod {\mathrm{E} _ {2}} \quad \text {(for all} \lambda \in \mathrm{K} \text {and all} x \in \mathrm{E} _ {1} \\ & \Leftrightarrow u (\mathrm{E} _ {1}) \subset \mathrm{E} _ {2} \end{array}
\]

whence the proposition.

COROLLARY 1. For \(\mathbf{E}_1\) to be stable under \(\rho\) , it is necessary and sufficient that \(\mathbf{E}_1\) be stable under \(\mathrm{L}(\rho)\) .

It suffices to take \(E_{1} = E_{2}\) in Proposition 13.

COROLLARY 2. Suppose that \(\rho\) is finite-dimensional. For \(\rho\) to be simple (resp. semi-simple), it is necessary and sufficient that \(\mathbf{L}(\rho)\) be simple (resp. semi-simple).

This follows from Corollary 1.

COROLLARY 3. Let \(x \in E\) . For \(x\) to be invariant under \(\rho(G)\) , it is necessary and sufficient that \(x\) be annihilated by \(L(\rho)(L(G))\) (that is that \(x\) be invariant under \(L(\rho)\) in the sense of Chapter I, § 3, Definition 3).

It suffices to take \(\mathbf{E}_1 = \mathbf{K}x\) and \(\mathbf{E}_2 = 0\) in Proposition 13.

COROLLARY 4. Let \(\rho'\) be another analytic linear representation of \(G\) on a complete normable space \(E'\) . Let \(T \in \mathcal{L}(E, E')\) . The following conditions are equivalent:

\begin{enumerate}
\def\labelenumi{(\roman{enumi})}
\item
  \(\mathrm{T}\rho (g) = \rho '(g)\mathrm{T}\) for all \(g\in \mathbf{G}\)
\item
  \(\mathrm{TL}(\rho)(a) = \mathrm{L}(\rho')(a)\mathrm{T}\) for all \(a\in \mathrm{L}(\mathrm{G})\)
\end{enumerate}

Let \(\sigma\) be the linear representation of \(G\) on \(\mathcal{L}(E, E')\) derived from \(\rho\) and \(\rho'\) ( \(\S 3\) , no. 11, Corollary 1 to Proposition 41). Condition (i) means that \(T\) is invariant under \(\sigma(G)\) . Condition (ii) means that \(T\) is annihilated by \(L(\sigma)(L(G))\) . It then suffices to apply Corollary 3.

COROLLARY 5. Suppose that \(\rho\) and \(\rho'\) are finite-dimensional. For \(\rho\) and \(\rho'\) to be equivalent, it is necessary and sufficient that \(\mathbf{L}(\rho)\) and \(\mathbf{L}(\rho')\) be equivalent.

This is a special case of Corollary 4.

COROLLARY 6. Suppose that \(\mathbf{G}\) is finite-dimensional. Let \(t \in \mathrm{U}(\mathrm{G})\) . For \(\mathbf{L}_t\) (resp. \(\mathbf{R}_t\) ) to be right (resp. left) invariant, it is necessary and sufficient that \(t\) belong to the centre of \(\mathrm{U}(\mathrm{G})\) .

For \(\mathrm{L}_t\) (resp. \(\mathbf{R}_t\) ) to be right (resp. left) invariant, it is necessary and sufficient that \(\varepsilon_g * t = t * \varepsilon_g\) for all \(g \in G\) , that is that \((\operatorname{Int} g)_* t = t\) . There exists an integer \(n\) such that \(t \in \mathrm{U}_n(G)\) . By Corollary 3 and Proposition 45 of § 3, no. 12, \((\operatorname{Int} g)_* t = t\) for all \(g \in G\) if and only if \([a, t] = 0\) for all \(a \in \mathrm{L}(G)\) , that is if and only if \(t\) commutes with \(\mathrm{U}(G)\) .

\subsection{NORMAL INTEGRAL SUBGROUPS}

Lemma 3. Let \(\mathbf{G}\) be a Lie group, \(\mathrm{H}_1\) and \(\mathrm{H}_2\) integral subgroups whose topology admits a countable base and \(g \in \mathbf{G}\) . Then

\[
g \mathrm{H} _ {1} g ^ {- 1} = \mathrm{H} _ {2} \Leftrightarrow (\operatorname{Ad} g) (\mathrm{L} (\mathrm{H} _ {1})) = \mathrm{L} (\mathrm{H} _ {2}).
\]

\(\operatorname{Ad} g = \mathrm{T}_{e}(\operatorname{Int} g)\) . Hence, by transport of structure, \((\operatorname{Int} g)(\mathrm{H}_{1})\) has Lie algebra \((\operatorname{Ad} g)(\mathrm{L}(\mathrm{H}_{1}))\) . As \(\mathrm{H}_{1}\) and \(\mathrm{H}_{2}\) have countable bases, to say that the sets \(\mathrm{H}_{2}\) and \((\operatorname{Int} g)(\mathrm{H}_{1})\) are equal amounts to saying that the integral subgroups \(\mathrm{H}_{2}\) and \((\operatorname{Int} g)(\mathrm{H}_{1})\) are equal (no. 2, Corollary 3 to Proposition 3). Then the lemma follows from Theorem 2 (i).

PROPOSITION 14. Let G be a Lie group and H an integral subgroup whose topology admits a countable base. The following conditions are equivalent:

\begin{enumerate}
\def\labelenumi{(\roman{enumi})}
\item
  H is normal in G;
\item
  L(H) is invariant under Ad(G).
\end{enumerate}

If further G is connected, these conditions are equivalent to the following:

\begin{enumerate}
\def\labelenumi{(\roman{enumi})}
\setcounter{enumi}{2}
\tightlist
\item
  L(H) is an ideal of L(G).
\end{enumerate}

If further \(\mathbf{G}\) is simply connected and \(\mathbf{L}(\mathbf{H})\) is of finite codimension in \(\mathbf{L}(\mathbf{G})\) , these conditions imply that \(\mathbf{H}\) is a Lie subgroup of \(\mathbf{G}\) and that \(\mathbf{G} / \mathbf{H}\) is simply connected.

The equivalence (i) \(\Leftrightarrow\) (ii) follows from Lemma 3. If further G is connected, condition (ii) is equivalent to saying that L(H) is stable under ad L(G) (no. 5, Corollary 1 to Proposition 13 and § 3, no. 12, Proposition 44).

Suppose that G is simply connected and that L(H) is an ideal of finite codimension in L(G). By Theorem 3 of no. 3, there exists a simply connected Lie group G' such that L(G') is isomorphic to L(G)/L(H). There exists a continuous morphism f of L(G) onto L(G') with kernel L(H). By Theorem 1 of no. 1, there exists a morphism \(\phi\) of G into G' such that L( \(\phi\) ) = f.~This morphism is a submersion and hence its kernel N is a Lie subgroup of G such that L(N) = Ker f = L(H). Hence H is the identity component of N and is therefore a Lie subgroup of G. Let \(\psi\) be the morphism of G/H into G' derived from \(\phi\) when passing to the quotient. This morphism is étale; since G is simply connected, \(\psi\) is an isomorphism of G/H onto G'.

COROLLARY 1. Let \(G\) be a finite-dimensional simply connected Lie group. Let \(m, h\) be Lie subalgebras of \(L(G)\) such that \(L(G)\) is the semi-direct product of \(m\) by \(h\) . Let \(M, H\) be the corresponding integral subalgebras of \(G\) . Then \(M\) and \(H\) are simply connected Lie subgroups of \(G\) and, as a Lie group, \(G\) is the semi-direct product of \(M\) by \(H\) .

By Proposition 14, H is a normal Lie subgroup of G and the Lie group G/H is simply connected. Let \(\pi\) be the canonical morphism of G onto G/H. There exists a morphism \(\theta\) of G/H into M such that L(θ) is the canonical isomorphism of L(G)/L(H) = L(G)/h onto L(M) = m. Then

\[
\mathrm{L} (\pi \circ \theta) = \mathrm{L} (\pi) \circ \mathrm{L} (\theta) = \mathrm{Id} _ {\mathrm{L} (\mathrm{G} / \mathrm{H})}
\]

and hence \(\pi \circ \theta = \mathrm{Id}_{\mathrm{G}/\mathrm{H}}\) . By no. 1, Corollary 1 to Proposition 1, \(\theta (\mathrm{G / H}) = \mathbf{M}\) . By Proposition 8 of § 1, no. 4, M is a Lie subgroup of G and the Lie group G is the semi-direct product of M by H.

COROLLARY 2. Let G be a finite-dimensional simply connected Lie group, H a normal connected Lie subgroup of G and \(\pi\) the canonical morphism of G onto G/H.

\begin{enumerate}
\def\labelenumi{(\roman{enumi})}
\item
  There exists an analytic mapping \(\rho\) of G/H into G such that \(\pi \circ \rho = \mathrm{Id}_{\mathrm{G / H}}\) .
\item
  For every mapping \(\rho\) with the properties of (i), the mapping \((h, m) \mapsto h\rho(m)\) of \(\mathbf{H} \times (\mathbf{G}/\mathbf{H})\) into \(\mathbf{G}\) is an isomorphism of analytic manifolds.
\item
  H and G/H are simply connected.
\end{enumerate}

Let \(n = \dim G - \dim H\) . The corollary is obvious for \(n = 0\) . We argue by induction on \(n\) .

Suppose that there exists an ideal of L(G) containing L(H) distinct from L(G) and L(H). Let H' be the corresponding connected Lie subgroup of G. Let \(\pi_{1}: G \to G/H'\) and \(\pi_{2}: H' \to H'/H\) be the canonical morphisms. By the induction hypothesis, there exist analytic mappings \(\rho_{1}: G/H' \to G\) , \(\rho_{2}: H'/H \to H'\) such that \(\pi_{1} \circ \rho_{1} = Id_{G/H'}\) , \(\pi_{2} \circ \rho_{2} = Id_{H'/H}\) . Let

\[
\pi_ {3} \colon \mathrm{G} / \mathrm{H} \rightarrow \mathrm{G} / \mathrm{H} ^ {\prime}
\]

be the canonical morphism. If \(x \in \mathrm{G} / \mathrm{H}\) and \(y\) is a representative of \(x\) in \(\mathrm{G}, y\) and \(\rho_1(\pi_3(x))\) have the same canonical image modulo \(\mathrm{H}'\) and hence \(x^{-1}\pi (\rho_1(\pi_3(x))) \in \mathrm{H}' / \mathrm{H}\) . Let

\[
\rho (x) = \rho_ {1} (\pi_ {3} (x)) \rho_ {2} (\pi (\rho_ {1} (\pi_ {3} (x))) ^ {- 1} x) \in G.
\]

Clearly \(\rho\) is an analytic mapping of G/H into G and

\[
\begin{array}{r l} \pi (\rho (x)) & = \pi (\rho_ {1} (\pi_ {3} (x))) \pi_ {2} (\rho_ {2} (\pi (\rho_ {1} (\pi_ {3} (x))) ^ {- 1} x)) \\ & = \pi (\rho_ {1} (\pi_ {3} (x))) \pi (\rho_ {1} (\pi_ {3} (x))) ^ {- 1} x = x. \end{array}
\]

If now the only ideals of \(\mathbf{L}(\mathbf{G})\) containing \(\mathbf{L}(\mathbf{H})\) are \(\mathbf{L}(\mathbf{G})\) and \(\mathbf{L}(\mathbf{H})\) , the Lie algebra \(\mathrm{L(G) / L(H)}\) is either 1-dimensional or simple. In both cases, \(\mathrm{L(G)}\) is the semi-direct product of a subalgebra by \(\mathrm{L(H)}\) ; this is obvious in the first case and in the second this follows from Chapter I, § 6, Corollary 3 to Theorem 5. Assertion (i) then follows from Corollary 1.

Assertion (ii) is obvious. Assertion (iii) follows from (i) and (ii).

The conclusions of Corollary 2 are no longer necessarily true when G is infinite-dimensional or when H is not normal (Exercises 8 and 15).

COROLLARY 3. Let \(\mathbf{G}\) be a finite-dimensional connected Lie group and \(\mathbf{H}\) a normal connected Lie subgroup of \(\mathbf{G}\) . The canonical morphism of \(\pi_1(\mathbf{H})\) into \(\pi_1(\mathbf{G})\) is injective.

Let \(G_{1}\) be the universal covering of G and \(\lambda\) the canonical mapping of \(G_{1}\) onto G. The Lie algebra of \(G_{1}\) is identified with L(G). The Lie subgroup \(\lambda^{-1}(H)\) of \(G_{1}\) is normal in \(G_{1}\) and its Lie algebra is L(H). Let \(H_{1}\) be the identity component of \(\lambda^{-1}(H)\) and \(\lambda_{1} = \lambda | H_{1}\) . Then \(\mathrm{L}(H_{1}) = \mathrm{L}(H)\) and hence \(\lambda_{1}\) is an étale morphism of \(H_{1}\) onto H. On the other hand, \(H_{1}\) is simply connected (Corollary 2) and hence is identified with the universal covering of H. Then, by General Topology, Chapter XI, the canonical morphism of \(\pi_{1}(H)\) into \(\pi_{1}(G)\) is identified with the canonical injection of Ker \(\lambda_{1}\) into Ker \(\lambda\) .

\subsection{PRIMITIVES OF DIFFERENTIAL FORMS WITH VALUES IN A LIE ALGEBRA}

PROPOSITION 15. Let \(\mathrm{G}\) be a Lie group, \(\mathrm{M}\) a manifold of class \(\mathrm{C}^r (r\geqslant 2)\) and \(\alpha\) a differential form of class \(\mathrm{C}^{r - 1}\) and degree 1 on \(\mathrm{M}\) with values in \(\mathrm{L}(\mathrm{G})\) , such that \(d\alpha +[\alpha ]^2 = 0\) . Suppose that \(\mathrm{M}\) is simply connected. For all \(x\in \mathrm{M}\) and all \(s\in \mathrm{G}\) , there exists one and only one mapping \(f\) of class \(\mathrm{C}^{r - 1}\) of \(\mathrm{M}\) into \(\mathrm{G}\) such that \(f(x) = s\) and \(f^{-1}.df = \alpha\) .

The uniqueness of \(f\) follows from § 3, no. 17, Corollary 2 to Proposition 59 and the fact that \(\mathbf{M}\) is connected. We prove the existence of \(f\) . There exist an open covering \((\mathrm{U}_i)_{i\in \mathbf{I}}\) of \(\mathbf{M}\) and, for all \(i\in \mathbf{I}\) , a mapping \(g_i: \mathrm{U}_i\to \mathrm{G}\) of class \(\mathrm{C}^{r - 1}\) such that \(g_i^{-1}.d g_i = \alpha\) on \(\mathrm{U}_i\) (§ 4, no. 6, Theorem 5). By § 3, no. 17, Corollary 2 to Proposition 59, \(g_i g_j^{-1}\) is locally constant on \(\mathrm{U}_i\cap \mathrm{U}_j\) . Let \(g_i g_j^{-1} = g_{ij}\) . Let \(\mathbf{G}_d\) be the group \(\mathbf{G}\) with the discrete topology. The \(g_{ij}: \mathrm{U}_i\cap \mathrm{U}_j\to \mathrm{G}_d\) are continuous and \(g_{ij}g_{jk} = g_{ik}\) on \(\mathrm{U}_i\cap \mathrm{U}_j\cap \mathrm{U}_k\) . Since \(\mathbf{M}\) is simply connected, there exist continuous mappings \(\lambda_i: \mathrm{U}_i\to \mathrm{G}_d\) such that \(g_i g_j^{-1} = \lambda_i\lambda_j^{-1}\) on \(\mathrm{U}_i\cap \mathrm{U}_j\) . Let \(g\) be the mapping of \(\mathbf{M}\) into \(\mathbf{G}\) whose restriction to \(\mathrm{U}_i\) is \(\lambda_i^{-1}g_i\) for all \(i\in \mathbf{I}\) . This mapping is of class \(\mathrm{C}^{r - 1}\) and \(g^{-1}d g = \alpha\) . The mapping \(f\) of \(\mathbf{M}\) into \(\mathbf{G}\) defined by \(f = s(g(x))^{-1}g\) satisfies the conditions \(f^{-1}.df = \alpha\) and \(f(x) = s\) .

\subsection{PASSAGE FROM LAWS OF INFINITESIMAL OPERATION TO LAWS OF OPERATION}

Lemma 4. Let \(\mathbf{G}\) be a connected topological group, \(\mathbf{X}\) a Hausdorff topological space and \(f_{1}, f_{2}\) laws of left (resp. right) operation of \(\mathbf{G}\) on \(\mathbf{X}\) such that, for all \(x \in \mathbf{X}\) , the mappings \(s \mapsto f_1(s, x)\) , \(s \mapsto f_2(s, x)\) of G into X are continuous. Suppose that there exists a neighbourhood V of \(\{e\} \times X\) in G \(\times\) X such that \(f_1\) and \(f_2\) coincide on V. Then \(f_1 = f_2\) .

Let \(x \in X\) and A be the set of \(g \in G\) such that \(f_{1}(g, x) = f_{2}(g, x)\) . Then A is closed in G. On the other hand, let \(g \in A\) ; we write \(y = f_{1}(g, x) = f_{2}(g, x)\) . There exists a neighbourhood U of e in G such that \(f_{1}(t, y) = f_{2}(t, y)\) for \(t \in U\) , in other words such that \(f_{1}(t', x) = f_{2}(t', x)\) for \(t' \in Ug\) (resp. gU). Hence A is open in G and therefore A = G.

PROPOSITION 16. Let G be a connected Lie group, X a Hausdorff manifold of class \(\mathbf{C}^r\) and \(f_1, f_2\) laws of left (resp. right) operation of class \(\mathbf{C}^r\) of G on X. If the laws of infinitesimal operation associated with \(f_1\) and \(f_2\) are equal, then \(f_1 = f_2\) .

By § 4, no. 7, Corollary to Proposition 11, there exists a neighbourhood V of \(\{e\} \times \mathbf{X}\) in \(\mathbf{G} \times \mathbf{X}\) such that \(f_{1}\) and \(f_{2}\) coincide on V. Hence \(f_{1} = f_{2}\) (Lemma 4).

Lemma 5. Let G be a simply connected topological group, X a Hausdorff topological space, U an open neighbourhood of e in G and \(\psi\) a continuous mapping of U × X into X such that \(\psi(e, x) = x\) and \(\psi(s, \psi(t, x)) = \psi(st, x)\) for all \(x \in X\) and s, t in U such that \(st \in U\) . Let W be a symmetric connected open neighbourhood of e such that \(W^{3} \subset U\) . There exists one and only one law of continuous left operation \(\psi'\) of G on X such that \(\psi'\) and \(\psi\) coincide on W × X. If G is a Lie group and X a manifold of class \(\mathbf{C}^r\) and \(\psi\) is of class \(\mathbf{C}^r\), then \(\psi'\) is of class \(\mathbf{C}^r\).

The uniqueness on \(\psi\) follows from Lemma 4. Let P be the permutation group of X. For \(u \in W^{3}\) , the mapping \(x \mapsto \psi(u, x)\) is an element \(f(u)\) of P and

\[
f \left(u _ {1} u _ {2} u _ {3}\right) = f \left(u _ {1}\right) f \left(u _ {2}\right) f \left(u _ {3}\right)
\]

for \(u_{1}, u_{2}, u_{3}\) in W. Applying Lemma 1 of no. 1, we obtain a morphism \(f'\) of G into P which extends \(f|\mathrm{W}\) . Let \(\psi'(g,x)=f'(g)(x)\) for all \((g,x)\in\mathbf{G}\times\mathbf{X}\) . Then \(\psi'\) is a law of left operation of G on X which coincides with \(\psi\) on W × X. As \(\psi'(g,\psi'(g',x))=\psi'(gg',x)\) for \((g,g',x)\in\mathbf{G}\times\mathbf{G}\times\mathbf{X}\) , the continuity of \(\psi\) on W × X implies the continuity of \(\psi'\) on gW × X for all \(g\in\mathbf{G}\) . Hence \(\psi'\) is continuous. If \(\psi\) is of class \(\mathbf{C}^r\), it is seen similarly that \(\psi'\) is of class \(\mathbf{C}^r\).

THEOREM 5. Let G be a simply connected Lie group, X a compact manifold of class \(\mathbf{C}^r\) ( \(r \geqslant 2\) ) and \(a \mapsto \mathrm{D}_a\) a law of left (resp. right) infinitesimal operation of class \(\mathbf{C}^{r-1}\) of L(G) on X. There exists one and only one law of left (resp. right) operation of class \(\mathbf{C}^{r-1}\) of G on X such that the associated law of infinitesimal operation is \(a \mapsto \mathrm{D}_a\) .

The uniqueness follows from Proposition 16. By § 4, no. 7, Corollary 1 to Theorem 6, there exist a neighbourhood V of \(\{e\} \times \mathbf{X}\) in \(\mathrm{G} \times \mathrm{X}\) and a law chunk of left (resp. right) operation of class \(\mathrm{C}^{r-1}\) of G on X, defined on V, such that the associated law of infinitesimal operation is \(a \mapsto \mathrm{D}_a\) . As X is compact, V can be assumed to be of the form \(U \times X\) , where U is an open neighbourhood of e in G. It then suffices to apply Lemma 5.
